\input{preamble}

\begin{document}

\title{Seminars in Algebraic Geometry}

\maketitle

\phantomsection

\label{section-phantom}

\tableofcontents



\section{Smoothable compactified Jacobians of
nodal curves}
\label{section-smoot-compa-jacob-nodal}

\noindent
Nicola Pagani, University of Liverpool and
Bologna.
Seminar of Algebraic Geometry UFF.
August 20, 2025.

\medskip\noindent
{\bf Abstract.} Building from examples, we
introduce an abstract notion of a
'compactified Jacobian' of a nodal curve. We
then define a compactified Jacobian
to be 'smoothable' whenever it arises as the
limit of Jacobians of smooth
curves. We give a complete combinatorial
characterization of smoothable
compactified Jacobians in terms of some 'vine
stability conditions', which we
will also introduce. This is a joint work
with Fava and Viviani.

\medskip\noindent



Let $C$ be a smooth curve and $d \in
\mathbb{Z}$. Define
$$
J^d_C=\{L:\text{$L$ is a line bundle of
degree $d$}\}/\sim
$$
which is a smooth projective variety of
dimension $g(C)$.

If $C$ is nodal we still can consider
$J^d_C$.
\begin{enumerate}
\item One connected component. Then the
Jacobian is $\mathbb{P}^1$ minus two
points. This is not universally closed, so it
is not proper.
\item Two components intersecting at one
point.
The pullback of the normalization splits the
degree in
intintely many ways, givieng that $J^{-1}_C$
is an infinite set of points. This
is not of finite type, so it is not proper.
\item The curve has two components
intersecting at two points.
This gives $J^{-2}_C$,
 which is a mixture of the two former items.
 (Probably not proper too.)
\end{enumerate}

\medskip\noindent
Now consider
$$
\text{TF}_C^d=\{\mathcal{F}:
\text{coherent on $C$, torsion-free, rank-1
on $C$}\}/\sim
$$
This satisfies the existence oart of the valu
point of properness.

Now we consider the moduli. Now we consider
the ideal sheaf of the (singular?)
point(s?):
\begin{enumerate}
\item One component. The stack is proper!
\item Two components intersecting once. Now
we get stacky points,
$x=[\bullet/\mathbb{G}_m]$. These points have
generic stabilizer. The resulting
stack is not separated because a morphism of
a curve, say $\mathbb{P}^1$ minus a
point … there are infinitely many ways to
extend a morphism from this thing to a
line bundle. So you cannot include any of
these stacky points. Recall that a
sheaf is {\it simple} if its automorphism
group is $\mathbb{G}_m$.
\item  The ideal sheaf of both nodes
$\mathcal{I}(N_1,N_2)$ has a positive
dimensional automorphism
group. The stack is not proper.
\end{enumerate}

\begin{definition}
\label{definition-fine-compactified-Jacobian}
A {\it fined compactified Jacobian} of $C$ is
an open connected substack of
$\text{TF}^d(C)$ that is also proper.
\end{definition}

\begin{remark}
\label{remark-algebraic-space}
This thing is automatically an algebraic
space.
\end{remark}

\begin{definition}
\label{definition-compactified-Jacobian}
A {\it compactified Jacobian} is an open
connected of $\text{TF}^d(C)$ that
admits a proper, good moduli space.
\end{definition}

Consider the Artin stack $\mathfrak{X}
\xrightarrow{\Gamma}X$ […]
 is a {\it good
moduli space} if
\begin{enumerate}
\item Every moduli factors
$$
\xymatrix{
\mathfrak{X} \ar[r]\ar[rd]&\mathcal{I}\text{
(ACC. space)}\\
& X
}
$$
\item
$\pi_*\mathcal{O}_{\mathfrak{X}}=\mathcal{O}
_X$.
\end{enumerate}

\medskip\noindent
We expect to find a notion of stability
condition to produce these things
[…] $[\bullet/ \mathbb{G}_m]$ would be the
polystable representative.

\begin{definition}
\label{definition-smoot-compa-jacob}
A compactified Jacobian $\overline{J_C}$ is
{\it smoothable} if all smoothings
$\mathcal{C}\to \Delta=\{0,\eta\}$ (with
$\mathcal{C}_0=C$),
$$
J^d_{\mathcal{C}_\eta}\cup C \to
\overline{J_C}
$$
is proper.
\end{definition}

\medskip\noindent

\begin{definition}
\label{definition-connected}
Let $X$ be a curve.
$$
\text{BCON}(X)=\{Y \subseteq X\text{s.t.
}Y,Y^c \text{ are connected}\}
$$
\end{definition}

\begin{definition}
\label{definition-v-satbility}
A $v$-curve is a generalization of items (2)
and (3) in the lists above
[it looks like two long snakes $\sim$ that
intersect several times, and $t$ is
the number of nodes].
A {\it $v$-condition} is a pair $n=(n_1,n_2)$
such that
$$
n_1+n_2=\begin{cases}
d+1-t\qquad &\text{we say the s.c. is
nondegenerate} \\
d-t\qquad &\text{degenerate}
\end{cases}
$$
$\mathcal{F}$ on $X$ is  {\it
$n$-(semi)stable} if
$\text{deg}\mathcal{F}_{X_i}>n_i$
($\text{deg}\mathcal{F}_{X_i}\geq n_i$)
for $i=1,2$.
\end{definition}
$\mathcal{F}_{X_i}=\mathcal{F}|_{X_i}$
torsion.
$$
\text{deg}(\mathcal{F}_{X_i})+\text{deg}
(\mathcal{F}_{X_2})
=d-|\text{sing}(F)|.
$$
Then
$$
\overline{J_C}(n)=\{\mathcal{F}\text{ is
semistable}\},
$$
a smooth compact Jacobian.

\begin{definition}
\label{definition-n-degeneracy}
A {\it degeneratation} of $v$-stab. on $X$ is
$n:\text{BCON}(X) \to \mathbb{Z}$
such that
\begin{enumerate}
\item
$$
n_Y+n_{Y^c}+|Y \cap Y^c|=
\begin{cases}
d+1\qquad &\text{ we say $Y$ is
$n$-nondegenerate} \\
d\qquad &\text{$Y$ is $n$-degenerate}
\end{cases}
$$
\item $Y_i$ no pa. common component
$n_{Y_1}+n_{Y_2}+\ldots$
$$
\xymatrix{
&Y_1\ar[dr]\ar[dl]\\
Y_2\ar[ur]\ar[rr]&&Y_3\ar[ll]\ar[ul]
}
$$
\end{enumerate}
\end{definition}

\begin{theorem}[-, et al]
\label{theorem-bijec-stabi-condi-nodal}
(bijection between stability conditions and
nodal curves) The map
$$
\left\{ \substack{\text{sm. comp.} \\
\text{Jac of $X$}} \right\} \to
\left\{ \substack{\text{$v$-stab.} \\
\text{cond. of $X$}} \right\}
$$
$$
n\mapsto
\overline{J_X}(n)=\{n\text{-semistable
sheaves}\}
$$
is a bijection. (The arrow should be from
right to left!)
\end{theorem}

F. Viviani had proved it for fine compact
Jacobians.

\section{Equivariant spaces of matrices of
constant rank}
\label{section-equiv-space-matri-const}

\noindent
Ada Boralevi, France.
Algebraic Geometry Seminar, IMPA.
August 27, 2025.

\medskip\noindent
{\bf Abstract.} A space of matrices of
constant rank is a vector subspace V, say
of dimension n+1, ofthe set of matrices of
size axb over a field k, such that
any nonzero element of V has fixed rank r. It
is a classical problem to look for
different ways to construct such spaces of
matrices. In this talk I will give an
introduction up to the state of the art of
the topic, and report on my latest
joint project with D. Faenzi and D. Fratila,
where we give a classification of
all spaces of matrices of constant corank one
associated to irreducible
representation of a reductive group.

\medskip\noindent

We are interested in vector spaces $U \subset
\text{Mat}_{m,n}(\mathbb{C})$,
with $m \leq n$, of {\it constant rank}, i.e.
such that for all $f \in U$,
$r:=\text{rank}f$ is the same.

Let $\ell(r,m,n):=\text{max}\dim U:U$ is of
rank $r$.

\medskip\noindent
{\bf Questions.}
\begin{enumerate}
\item $\ell(r,m,n)=$? In general not known.
\item Find relations among $\ell, r,m$ and
$n$.
\item Construction of examples and
classification.
\end{enumerate}

\begin{example}
\label{example-what-we-know}
\begin{enumerate}
\item $\ell(1,m,n)=n$,
$$
\begin{pmatrix}
x_1&x_2&\cdots &x_n\\
\\
\\
\\
\end{pmatrix}
$$
\item $\text{rank}=2$? There are two cases
([Atkinson '83], [Eisenbud-Haus '88])
\begin{itemize}
\item Compression spaces,
$$
\begin{pmatrix}
*&*&\cdots &*\\
*&0&\cdots & 0\\
\vdots &\\
*&0&\cdots&0
\end{pmatrix}
$$
\item Skew-symmetric matrices of $3 \times
3$.
\end{itemize}
\item $\ell(r,m,n) \geq n-r+1$. Because you
can put a matrix of
$m\times n$ (with the first $r$ rows that can
have nonzero entries):
$$
\begin{pmatrix}
x_1 &  x_2 &  \cdots & & x_{n-r+1}\\
& x_1 & x_2 & \cdots &\\
\\
\\
\\
\end{pmatrix}
$$
\end{enumerate}
\end{example}

\begin{theorem}[Westwick '86]
\label{theorem-Westwick}
\begin{enumerate}
\item $n-r+1 \leq \ell(r,m,n) \leq 2n-2r+1$.
\item If $n-r+1 \not|
\frac{(m-1)!}{(r-1)!}\implies \ell$
\end{enumerate}
\end{theorem}

We can see these spaces as (subvarieties?) of
determinantal varieties
$M_n=\{f \in
\text{Mat}_{m,n}(\mathbb{C}):\text{rank}(f)
\leq  r\}$.
[Interpretation via secant varieties inside
Segre embedding.]

\medskip\noindent
Consider a map $\varphi: U \to \Hom(V,W)$.
Then $\varphi \in U^* \otimes  V^* \otimes
W$.
We get that
$\varphi$ is of constant rank if and only if
some kernel, image and cokernel are vector
bundles.

\medskip\noindent
{\bf Focus of today.} What happens when
$U,V,W$ are irreducible
representations of a complex reductive group
$G$?

\medskip\noindent
{\bf Question.} What is the natural
equivariant morphism
$$
U \to \Hom(V,W)=V^* \otimes  W
$$
of constant rank?

\medskip\noindent
Consider the case of $G=\text{SL}_2$. All the
irreducible representations (which are
self-dual) are given by
$V(m-1)\cong
\mathbb{C}[x,y]_{\text{deg}=m-1}$.

Recall the Clebsh-Gordon decomposition ($m
\leq  n$)
$$
V(m-1)\otimes
V(n-1)=\bigoplus_{i=0}^{m-1}V(n-m+2i)
$$

\begin{theorem}[B-Faenzi-Lella '22]
\label{theorem-BFL22}
$$
V(n+m-2)\hookrightarrow \Hom(V(m-1),V(n-1))
$$
is of constant rank (corank 1) if and only if
$$
n-m+2i|m-1
$$
\end{theorem}

\begin{theorem}[B. Faenzi, Fratila '25]
\label{theorem-BFF25}
Let $V(\nu)$, $V(\mu)$ and $V(\lambda)$ be
irreducible representations
of a comple reductive group $G$, with
$$
\dim(V(\mu)) \leq \dim(V(\lambda))_-
$$
If there exists a morphism of representations
$$
\varphi: V (\nu) to \Hom (V (\lambda),V(\mu))
$$
then $\varphi$ is of constant corank 1 if and
only if
there exists a simple root $\alpha_i$ such
that
\begin{enumerate}
\item $\lambda=\mu+\nu-\alpha_i$,
\item $\nu$ is a multiple of $\nu$,
\item  $\nu$ is a multple of $\omega_i$
\end{enumerate}
\end{theorem}

\section{Irreducibility of the Hilbert scheme
of points and the class of 2-step
ideals}
\label{section-irred-hilbe-schem-point}

\noindent
Michele Graffeo, SISSA.
Algebraic Geometry Seminar, UFF.
September 29, 2025.

\medskip\noindent
{\bf Abstract.}
Hilbert schemes of points on a
quasi-projective variety X are classical
objects
in algebraic geometry. Roughly speaking, they
parametrise ideals of a polynomial
ring with complex coefficients having finite
colength. Although Hilbert schemes
always have a distinguished component called
the smoothable component, their
geometry is quite pathological and one of the
main open problems around them
concerns their irreducibility. In a joint
work with Giovenzana, Giovenzana,
Lella we introduce the notion of 2-step
ideals. We show that the loci
parametrising these ideals are in general not
contained in the smoothable
component of the Hilbert scheme, thus
providing new examples of
extra-components.  In my seminar I will
discuss our class of ideals and relate
them to the compressed algebras considered by
Iarrobbino in the eighties.
Finally, I will show how to extend our result
to the nested setting.


\medskip\noindent

As usual, you start with a functor and
prove representativity.

\begin{definition}
\label{definition-hilbert-scheme}
Let $X$ be a smooth
quasi-projective variety and a nonegative
integer $d$.
Define
$$
\begin{aligned}
\underline{Hilb}^d: \text{Sch}_\mathbb{C}
&\longrightarrow \mathsf{Set} \\
B &\longmapsto \left\{Z \hookrightarrow B
\times X:
\substack{B\text{-flat} \\ B\text{-finite}\\
\text{length }d}\right\}
\end{aligned}
$$
\end{definition}

\begin{theorem}[Grothendieck]
\label{theorem-grothendieck}
$\underline{Hilb}^d(X)$ is represented by a
quasi-projective
scheme $Hilb^d(X)$
\end{theorem}

Idea: $\mathbb{C}$-points of $Hilb^d(X)$
are in correspondence with $Z
\overset{\subset}{\text{finite}}X$
of length.

A fat point $Z= \text{Spec}(A)$ $(A,
\mathfrak{m})$ is a
local artinian $\mathbb{C}$-algebra of finite
type.

If $X$ smooth can consider $X=\mathbb{A}^n$,
and it's the same to consider
\begin{itemize}
\item $Z \hookrightarrow \mathbb{A}^n$ of
length $d$.
\item $I \subset \mathbb{C}[x_1,\ldots,x_n]$
of
colength $d$.
\item $A= \mathbb{C}[x_1,\ldots,x_n]/I$ with
$\dim_\mathbb{C}=d$.
\end{itemize}

\begin{definition}
\label{definition-nested-Hilbert-scheme}
$X$, $0 \leq d_1 \leq \ldots
d_r=\underline{d}$.
$$
Hilb^{\underline{d}}(X)=
\{Z_1\not\hookrightarrow
\ldots
\not\hookrightarrow Z_r  \not\hookrightarrow
X:
\text{len}(Z_i)=d_i\}
$$
\end{definition}

\begin{theorem}[Hartshorne (r=0), Fogarty,
Kalpan (r>1)]
\label{theorem-HK}
$Hilb^d(X)$ is connected.
\end{theorem}

For $r=1$, $Hilb^d(\mathbb{A}^n)$ is smooth
iff
$n \leq 2$, $d \leq 3$.
If $r>1$, $n=2$, it is smooth iff $r=2$ and
$d_1=d_2=-1$.
If $r>1$ and $n>2$, it is smooth iff $r=2$,
$(d_1,d_2)=\{(1,2),(2,3)\}$.

Iror components.
\begin{itemize}
\item $(r=1)$ $n\geq 4$,
$Hilb^d(\mathbb{A}^3)$
is irreducible iff $d\leq 7$ ($\impliedby$
[E-I], $\implies $ [Mozzola].
$n=3$ irredicuble if $d \leq  11$ (8,9 Sivic;
10 Jardim et al; 11, Jelisievv et
al.)
\item $(r=2)$, $n=2$, irreducible
[G-Rosul-Sebastian].
\item $(r \geq 3)$ $n=2$, there exist $d_1
\leq \ldots \leq d_5$
such that
$Hilb^{\underline{d}}(\mathbb{A}^i)$ is
irreducible [S-R].
\end{itemize}

There is an analogy between classical in
dimension 3
and nested in dimension 2.

Schematic structure.
\begin{itemize}
\item [I] $Hilb^{21}(\mathbb{A}^4)$ has at
least a generically
non-reduced component.
\item Problem: what about $n=3$?
\item Theorem. If $Hilb^{d_1,\ldots
d_r}(\mathbb{A}^i)$ is
irreducible $\implies $
$Hilb^{(1,d_1,\ldots,d_r)}(\mathbb{A}^n$
has a generically non reduced component.
\end{itemize}

Other results.
\begin{itemize}
\item …
\end{itemize}
Hilbert scheme function.
For a local artinian $\mathbb{C}$-algebra of
finite type there is an associated graded
$\text{Gr}_\mathfrak{m}(A)
=\bigoplus_{i \geq
0}\mathfrak{m}^1/\mathfrak{m}^{i+1}$
which allows us to define the Hilbert
function
$$
\begin{aligned}
h_A: \mathbb{Z} &\longrightarrow \mathbb{N}
\\
i &\longmapsto \dim_\mathbb{C}
\mathfrak{m}_i/\mathfrak{m}_{i+1}
\end{aligned}
$$
sistability for graded $A$-modules.
The socol is the annihilator of
the ideal which gives the algebra
e.g. $\mathbb{C}[x,y](x^2,xy,y^2)$ then
$\text{Soc}(A)=\left\langle{}x,y^2\right\rangle{}_\mathbb{C}
$.

\medskip\noindent
\begin{definition}
\label{definition-smoothable-component}
The {\it smoothable component} is
 $$
V_{sm}=\overline{\{[(z_i)_{i=1}^n \in
Hilb^{\underline{d}}(\mathbb{A}^n):
Z_r\text{ is reduced}\}}
$$
\end{definition}

\begin{definition}
\label{definition-elementary-componenet}
$V \subset
Hilb^{\underline{d}}(\mathbb{A}^n)$ is
an {\it elementary component} if
$\forall [Z_1,\ldots,Z_r] \in V$
then $Z_r$ is a fat point.
\end{definition}

Instead of looking for all components of the
Hilbert scheme
we look for elementary components. This may
not be some simple;
sometimes there's infinitely many. There
is no method to find them.

\begin{theorem}[Irrabaro]
\label{theorem-I}
Every irreducible component $V \subset
Hilb^d(\mathbb{A}^n)$
is generically étale locally product of
elementary components.
\end{theorem}

\begin{definition}
\label{definition-}
Let $h:\mathbb{Z} \to \mathbb{N}$ be a
function
with finite support, $|h|=\sum_{ i \in
\mathbb{Z}}h(i)=d$. Define
\begin{itemize}
\item $H_n=\{[A] \in Hilb^d(\mathbb{A}^n):
h_A \equiv h\}$.
This is closed by semicontinuity of the
function; but can be expressed
a by a representable functor.
\item The following also has a canonical
schematic structure
$$
\begin{aligned}
\pi_n:H_n  &\longrightarrow
\mathcal{H}_n=\{[A] \in H_n: A\text{ is
graded}\} \\
A &\longmapsto A=\text{Gr}_{\mathfrak{m}}(A)
\end{aligned}
$$
Where the first is $\not\hookrightarrow $ in
the second.
\end{itemize}
\end{definition}

Recall that
$$
T_{[I]}Hilb^d(\mathbb{A}^n)=\Hom_R(I,R/I)
$$
\begin{theorem}[B-B,J,G G G L]
\label{theorem-B-B-J-G-G-G-L}
There is a graded $[I] \in \mathcal{H}_n$.
Then
$$
\begin{aligned}
T_{[I]}\mathcal{H}_n&=\Hom_R(I,R/I)_{=0}\\
T_{[I]}\pi_n([I])&='>0\\
T_{[I]}H_n&=' \geq 0.
\end{aligned}
$$
\end{theorem}

\begin{definition}
\label{definition-compressed}
$(A,\mathfrak{m}_A)$ is {\it compressed} is
for all  $(A',\mathfrak{m}')$
with $e(A)=e(A')$ and $len(A) \geq  len(A')
$.
\end{definition}

\begin{theorem}
\label{theorem-E}
$E_h \subset H_n$ is $H$-compressed
$E_h \neq  \emptyset$
\end{theorem}

\section{Normal and Restricted Tangent
Bundles of Rational
Curves in Hypersurfaces}
\label{section-restricted-tangent-bundles}

\noindent
Lucas Mioranci, IMPA.
Algebra Seminar, IMPA.
October 8, 2025.

\medskip\noindent
{\bf Abstract.}
Let $X\subset \mathbb{P}^n$ be a degree $d$
hypersurface containing a smooth
rational curve $C$ of degree $e$. The normal
bundle $N_{C/X}$ and the restricted
tangent bundle $T_X|_C$ split as direct sums
of line bundles of the form
$\bigoplus_i
\mathcal{O}_{\mathbb{P}^1}(a_i)\qquad$ called
their splitting type.
The splitting type of $N_{C/X}$ and
${T_X}|_C\;\;$ controls the local structure
of the space of rational curves and
determines how many general points of $X$ we
can interpolate by deforming the curve $C$.
By combining explicit computations
of $T_X|_C\;\;$ and an induction argument, I
classify all triples $(e,d,n)$ such
that a general degree $d$ hypersurface
$X\subset \mathbb{P}^n$ contains a
rational curve $C$ of degree $e$ whose
restricted tangent bundle ${T_X}|_C\;\;$
is balanced. The case of quadrics is
particular, in which we show that
odd-degree rational curves do not interpolate
the expected number of points on
quadric hypersurfaces.  For the normal
bundle, I compute explicit examples of
hypersurfaces $X$ for all possible splitting
types of $N_{C/X}$ when $C$ is the
rational normal curve. Additionally, for
$d\ge 3$, we compute the dimension of
the space of hypersurfaces $X$ such that
$N_{C/X}$ has a given splitting type.



\medskip\noindent

$\mathcal{O}$ with no subindex means
$\mathcal{O}_{\mathbb{P}^1}$.
$$
\xymatrix{
&0\ar[d]&0\ar[d]\\
&T_{\mathbb{P}^1}=\mathcal{O}_(2)\ar[d]\ar[r]
^=
&T_{\mathbb{P}^1}=\mathcal{O}(2)\ar[d]\\
0\ar[r]&T_{X}|_{C}\ar[d]\ar[r]&T_{\mathbb{P}
^n}|_{C}\ar[d]\ar[r]&
N_{X/\mathbb{P}^n}|_{C}\cong
\mathcal{O}(de)\ar[d]_=\ar[r]&0\\
0\ar[r]&
N_{C/X}\ar[d]\ar[r]&N_{C/\mathbb{P}^n}\ar[d]
\ar[r]&
N_{X/\mathbb{P}^n}|_{C}\cong
\mathcal{O}(de)\ar[r]&0\\
&0&0
}
$$
Throughout this talk
$X$ is a degree $d$ hypersurface in
$\mathbb{P}^n$
and $C \subset X$ is a smooth rational curve
of degree $e$.

Since $C$ is rational we can think of it
as a map $f:\mathbb{P}^1 \to X$.
Then we can pullback any bundle over $C$
back to $\mathbb{P}^1$,
and by Birkhoff-? Theorem
we know that vector bundles over
$\mathbb{P}^1$
split as
$$
E\cong \bigoplus_{i=1}^r
\mathcal{O}_{\mathbb{P}^1}(a_i),
\qquad a_1\leq \ldots \leq a_n.
$$
This decomposition is called the
{\it splitting type} of $E$.
We say $E$ is {\it balanced} if
$|a_i-a_j|\leq 1$ for all $i,j$.

Balancedness is an open condition,
i.e. if you have a family of bundles over
$\mathbb{P}^1$
and one of them is balanced, then all of them
are,
(also you can think it's a generic
condition).

We know the normal bundle $N_{C/X}$ has rank
$n-2$
and degree  $e(n-d+1)-2$.
And the tangent bundle $T_{X/C}$
has rank $n-1$ and degree $e(n-d+1)$.

\begin{example}
\label{example-conic}
A conic $C \subset \mathbb{P}^3$.
Since a conic in $\mathbb{P}^3$ is contained
in a plane,
this obliges to have a ``distinguished normal
direction'',
$\mathcal{O}(2)$. Indeed,
$$
\begin{aligned}
N_{C/\mathbb{P}^3}&\cong(N_{Q/\mathbb{P}^3}
\oplus
N_{\mathbb{P}^2/\mathbb{P}^3}|_{C}\\
&\cong(\mathcal{O}(2) \oplus
\mathcal{O}(1))|_{C}\\
&\cong\mathcal{O}(1)\oplus \mathcal{O}(2).
\end{aligned}
$$
\end{example}

\medskip\noindent
Now we discuss interpolation of general
points.
Let $p_1,\ldots,p_n \in \mathbb{P}^1$
and $x_1,\ldots, x_n \in X$
be general points.
Let $f:\mathbb{P}^1 \to X$ be such that
$f(p_i)=x_i$.

Consider the space of morphisms mapping
the  $p_i$ to $x_i$. It's known that its
tangent space
at $f$
is given by
$$
T_{[f]}\Mor(\mathbb{P}^1,X,p_i \mapsto x_i)
\cong H^0(\mathbb{P}^1,f^*T_X(-n)).
$$
\noindent
Deformations of $f$ (mapping $p_i \mapsto
x_i$)
dominate $X$ if $a_i-n_i \geq 0$.
Also, deformations of $f$ interpolate up
to $a_1+f$ general points in $X$.

\medskip\noindent
How many points can deformations of the curve
interpolate?
This means (I think) that the deformations of
the curve
are still such that $p_i \mapsto  x_i$.
``And for the normal bundle, interpolations
means that the curves contain the points''.

If $T_X|_{C}$ is balanced, $C$ can
interpolate up to
$$
\left\lfloor \frac{e(n+1-d)}{n_1}
\right\rfloor +1
$$
general points. The case of $N_{C/X}$ is
similar.

\medskip\noindent
Next we survey some relevant results.

\begin{theorem}[Coskun-Riedl, 2018]
\label{theorem-CR18}
A general Fano $X$ of degree $d \geq 2$
contains a degree $c$ rational smooth curve
 $C$ with balanced $N_{C/X}$
for every $1 \leq e \leq  n$.
\end{theorem}

The idea is to observe that the space of
curves has
balanced normal bundle.

\begin{theorem}[Ran, 2021]
\label{theorem-R21}
Extends the above for
$1 \leq  e \leq 2n-2$
and $d \geq 4$.
\end{theorem}

Idea: general hypersurface contains a curve
of balanced normal bundle.

\begin{theorem}[Ran, 2024]
\label{theorem-R24}
$X$ general Fano hypersurface.
There exist smooth rational curve with  $C$
with balanced tangent bundle $T_{X}|_{C}$
for $e$ in some arithmetic progressions.
\end{theorem}

$e$ is in general very large.
They show there are curves of arbitrarily
large
degrees with balanced tangent bundle.

\medskip\noindent
Now we describe the new results.

\begin{theorem}
\label{theorem-Lucas}
Let $X \subseteq \mathbb{P}^n$ be Fano
hypersurface of degree $3 \leq d \leq n$.

\begin{enumerate}
\item If $C \subset X$ is a rational smooth
curve
of degree $e$ with
$e \leq  \frac{n-1}{n+1-d}$,
then $T_X|_{C}$ is not balanced.

\item A general hypersurface $X$
contains rational curves $C$
of degree $e$ with balanced
$T_X|_{C}$ for every
$e> \frac{n-1}{n+1-d}$.
\end{enumerate}
\end{theorem}

\begin{theorem}
\label{theorem-lucas2}
Let $X \subseteq \mathbb{P}^n$ be a smooth
quadric hypersurface.

\begin{enumerate}
\item If $e \leq 2$ is even,
then exists a curve $C$ of degree $e$
with $T_X|_{C}\cong \mathcal{O}(e)^{n-1}$.
(I.e. we can interpolate $e+2$ points
if $e$ is even.)

\item If $e \geq 1$ is odd,
then there is no curve $C$ of degree $e$
with balanced $T_X|_C$.
The best we can do is the most balanced
bundle
before balancedness, namely
$T_X|_C \cong
\mathcal{O}(e-1)\oplus\mathcal{O}(e)
\oplus\mathcal{O}(e+1)$.
(I.e. we can only interpolate $e$ points if
$e$ is odd.)
\end{enumerate}
\end{theorem}

Next we give the main ideas in the proof.
Let $C$ be the rational normal curve
of degree $e$ in $\mathbb{P}^1$,
$$
\begin{aligned}
C:\mathbb{P}^1  &\longrightarrow \mathbb{P}^n
\\
(s:t) &\longmapsto
(s^2:s^{e-1}t:\ldots:t^e:0:\ldots:0).
\end{aligned}
$$
$X=V(F)$.

Then find some examples. Then,

\begin{proposition}
\label{proposition-proposition}
If $\mu(T_X|_{C})=\frac{e(n+1-d)}{n-1}\leq 1$
then $T_X|_{C}$ is not balanced.
\end{proposition}

\noindent
Use the proposition to show that $T_X|_C$
is not balanced if $d \leq  n$
and $e le \frac{ n-1}{n+1-d}$.

This implies Theorem \ref{theorem-Lucas}.

\begin{lemma}
\label{lemma-split-normal-bundle}
If $N_{C/X}\cong \bigoplus_i
\mathcal{O}(a_i)$
with $a_i < 4$ for all $i$,
then $T_{X}|_C \cong N_{C/X} \oplus
\mathcal{O}(2)$.
\end{lemma}

\noindent
This implies that
$\text{Ext}^1(N_{C/X},\mathcal{O}(2))=0$.

The most important step is the following,
used as the induction step:

\begin{proposition}
\label{proposition-induction}
If $T_Y|_C$ is balanced for some degree $d$
hypersurface $Y \subseteq \mathbb{P}^{n-1}$,
then there exists a degree $d$ hypersurface
$C \subseteq \mathbb{P}^n$ with balanced
$T_X|_C$.
\end{proposition}

\noindent
Combining the examples, Lemma
\ref{lemma-split-normal-bundle}
and Proposition \ref{proposition-induction},
we establish that for $e \leq
\text{max}\{2d-2,n\}$.

The following argument is used to show that
actually what we showed so far is enough for
all degrees.

\begin{lemma}
\label{lemma-smoothing-argument}
$C=C_1 \cup  C_2$, $C_1 \cong \mathbb{P}^1$.
$E$ vector bundle on $C$ such that
$E|_{C_1}$ is balanced,
$E_{C_2}$ is perfectly balanced
(numerical condition that works when
$e=n-1$).
Then if $E$ is the specialization of a family
of vector bundles $E'$ on $\mathbb{P}^1$,
then $E'$ is balanced.
\end{lemma}

\noindent
Basically, glue and smooth low degree curves
(i.e. $e \leq  \text{max}\{2d-2,n\}$)
to get all degrees


\medskip\noindent
Now we prove the result for quadrics,
i.e. $X$ is a quadraic.
For even $e$ compute examples and conclude
 $T_X|_C \cong \mathcal{O}(e)^{n-1}$.

For odd $e$, compute examples,
$T_X|_C\cong \mathcal{O}(c-1) \oplus
\mathcal{O}(2)^{3}\oplus \mathcal{O}(-1)$.

\begin{proposition}
\label{proposition-interpolation2}
$n\geq 3$, $m\leq  e+1$.
There is a degree $e$ curve $\phi_e$
interpolating
$m$ general points if and only if there is a
degree $e-2$ curve $\psi_{e-2}$ interpolating
$m-2$ general points in $X$.
\end{proposition}

\noindent
The idea is to do induction.
Suppose you can interpolate $e$ points,
then show you can interpolate $e+2$.

Chose $y_1,\ldots,y_{m-2}$ on the quadric.
Choose a line in $\mathbb{P}^{m-1}$
intersecting the quadric,
and $x_1,\ldots,x_{m-1}$ points on the line.
Join the points on the line to the points on
the quadric,
Then you get somem points of intersection
with the quadric.
Those points and the two points of
intersection of the line and the quadric
let us construct a scroll, which intersects
the quadric.
The intersection of the scroll and the
quadric
(I guess the scroll contains the lines, and
so the points $y_i$)
gives the curve passing through the $y_i$.

\section{Regular but non smooth curves of
genus 3}
\label{section-regular-but-non-smooth}

\noindent
Cesar Hilario, .
Algebra Seminar, IMPA.
October 15, 2025.

\medskip\noindent
{\bf Abstract.} Regular but non-smooth curves
are a unique feature of
geometry in positive characteristic, that
results from the fact that over an
imperfect field the notion of regularity is
weaker than the notion of
smoothness. In the setting of algebraic
geometry over an algebraically closed
field, these curves correspond to fibrations
by singular curves, which are
fibrations of relative dimension 1 whose
fibers are singular. The most famous
examples are arguably the so-called
quasi-elliptic curves and quasi-elliptic
fibrations, which play a key role in the
Bombieri-Mumford classification of
algebraic surfaces in characteristics 2 and
3. In this talk I will discuss the
case of genus 3 in characteristic 2, with an
eye towards the classification of
regular plane projective quartic curves that
become rational after base change.

\medskip\noindent
arXiv 2409.05464.

\medskip\noindent
$C$ will denote a projective geometrically
integral curve over
a not necessaily closed field $K$.
Geometricallt integral
means that $C \otimes_K \overline{K}$ is
integra.
Suppose that $p=\text{char}(K) \geq 0$
and $g=h_1(\mathcal{O}_C)$ is of genus $C$.

We define
\begin{enumerate}
\item $C$ regular iff local rings of $C$ are
regular (DVR).
\item $C$ smooth iff $C$ regular and $K
\otimes_K \overline{K}$ regular.
\end{enumerate}

\noindent
Thus it's obvious that smoothness implies
regularity.
In this talk we explain why the converse is
not true.

Note that if $K$ is perfect (e.g. $p=0$ or
$K=\overline{K}$) implies
 smooth = regular. Recall that $K$ is
 imperfect iff $p>0$
the powers of $K$, $K^p$, is a proper
subfield of $K$
(it's always a subfield but we ask it's
proper).

\medskip\noindent
Assume that $C$ is regular and $p>0$.
The genus of rthe normalization is less than
or equal to the genus of the curve,
i.e.
$$
\tilde{g}=h^1(\mathcal{O}_{\widehat{C
\otimes_K \overline{K}}})\leq g.
$$
We know that $C$ is smooth iff
$\overline{g}=g$.
So $g-\overline{g}$ is a measure of
non-smoothness.

\begin{theorem}[Tate]
\label{theorem-tate}
$\frac{p-1}{2}$ divides $g-\overline{g}$.
\end{theorem}

As a corollary, we get that if $C$ is
non-smooth,
$g-\overline{g}>0$ and so $p \leq  2g+1$.

\begin{example}
\label{example-attain}
There are cases in which the bound is sharp,
namely $C:y^2=x^p+t$, with $t \in K\setminus
K^p$.
Then $g=\frac{p-1}{2}$ and $\bar{g}=0$.
Over $\overline{K}$ rhere exists a singular
point,
$(x,y)=(-t^{1/p},0)$ which
{\bf is not visible over $K$}, that is we get
that
$t^{1/p}\not \in K$.
\end{example}

\begin{example}
\label{example-elliptic}
$\{\text{elliptic curves}\}=\{g =
\overline{g}=1,\text{$\exists K$-rational
point }\}=\{\text{smooth cubic curves with
$K$-rational point}\}$.
$y^2=x^3+4g_i x+27$.
\end{example}

\begin{example}
\label{example-quasielliptic}
$g=1$, $\overline{g}=0$, $p=2,3$,
quasielliptic curves. (Important in the
classification of surfaces
in positive characteristic.)
\end{example}

\medskip\noindent
Now we make a digression to describe
Kodaira-Enriques (over $\mathbb{C}$)
in positive characteristic. Bombieri-Mumford:
{\bf new objects} appear,
quasielliptic surfaces.

Let $k=\overline{k}$, $S$ smooth elliptic
({\it quasielliptic}) surface over $k$.
Suppose we have a fibration
where the general fiber is an elliptic curve
an elliptic curve (resp. a plane
cuspidal cubic) and the generic fiber an
elliptic (resp. quasielliptic)
curve over $k(B)$.

$$
\xymatrix{
S\ar[d]\\
B
}
$$

\medskip\noindent
Back to our main content, let's assume that
$g=3$.
$$
\xymatrix{
\{g=1,\exists K\text{-rat. pt}\}\ar[r]&\ar[l]
\left\{\substack{\text{reg. cubic curves} \\
\text{with $K$-rat pt}}\right\}
}
$$
Recall that genus-degree formula
$$
g=\frac{(d-1)(d-2)}{2}=1,
$$
So then $d=3,4$ gives  $g=3$ (right?).

No we present our main theorem, which is a
classification result:

\begin{theorem}
\label{theorem-classification-quartics}
\begin{itemize}
\item $C$ regular non-hyperelliptic curve
over $K$.
\item  $g=3$, $\overline{g}=0$.
\item $p=2$.
\end{itemize}
(Not that non-hyperllipticity of $C$ and
$g=3$ imply that $C$ is a quartic
over $K$, and that  $\overline{g}=0$ implies
$C$ is geometrically rational.)
Then $C$ is isomorphic to one of the
following quartics.

\begin{enumerate}
\item $y^4+az^4+xz^3+bx^2z^2+cx^4=0$ where
$a,b,c \in K$, $c \not \in K^2$.
(Notice the imperfectness of $K$ is crucial.)

\item $y^4+az^4+bx^2y^2+cx^2z^2+bx^3z+dx^4=0$
for $a,b,c,d, \in K$,
$a \not \in K^2$, $b\not =0$.

\item
\item
\item
\end{enumerate}
\end{theorem}

\noindent
Now we enumerate properties of these
families.

\begin{itemize}
\item $C$ is a purely inseparable double
cover of a quasielliptic curve.
To understand this consider the induced map
from
Frobenius map $K \to K$, $a \mapsto a^p$ and
do base change:
$$
\xymatrix{
C^{(p)}\ar[r]\ar[d]\ar@{}[dr]|-{\lrcorner}&
C\ar[d]\\
\Spec K\ar[r]^{\substack{\text{induced from}
\\ \text{Frobenius}}}&\Spec K
}
$$
This gives a map $C \to C^{(p)}$ by universal
property of pullback.

$$
\xymatrix{
C \ar[rrrd]^{\text{Frobenius}}
\ar@{-->}[rrd]\ar[rrdd] & & \\
& &  C^{(p)}\ar[d] \ar[r] & C \ar[d] \\
& &  \Spec K\ar[r] & \Spec K
}
$$

Then consider the normalization of
$C^{(p)}=C_1$ which
turns out to be quasielliptic; by universal
property of normalization we can lift:
$$
\xymatrix{
& \widehat{C^{(p)}}\ar[d]\\
C\ar@{.>}[ur]\ar[r]&C^{(p)}.
}
$$
Here the lower arrow is purely inseparable of
degree 2.

\item There exists a unique non smooth point
$x \in C$.

\item Now we iterate the construction of the
first item:
$$
\xymatrix{
\underbrace{C}_{g=3}\ar[r]&\underbrace{C_1}
_{g=1}\ar[r]&
\underbrace{C_2}_{g_2-9}\ar[r]&
\underbrace{C_3}_{g_3=0}\ar[r]&
\underbrace{C_4}_{g_4=0}\ar[r]&\cdots
}
$$
Let $x$ be the non-smooth (non-rational)
point of $C$. Map it under
these maps to $x_1$( non-smooth,rational),
then  $x_2$ (smooth, rational?),
then $x_3 $ (smooth, rational.), then $x_4$
(smooth, rational), etc.

We have
$$
\begin{matrix}
&\text{$x_2$ rational}& \text{$x$ canoical
divisor} &
E=K(c_2)\\
(i)& Y & Y & Y\\
(ii) & N & Y & N\\
(iii) & Y & N & N\\
(iv) & N & N & Y\\
(v) & N & N & N
\end{matrix}
$$
Where the last column is explained next.

\item Canonical field of $C$ = subfield of
$K(C)$ generated
by the quotients of all non-zero holomorphic
differentials
=  $K(C)$ ($C$ non-hyperelliptic).

\item Pseudocanonical field of $C$ = …
non-zero \underline{exact}
holomorphic differential := $E$. So one has
$[K(c),E]=4=p^2$.

$$
\xymatrix{
&K(C)\ar[ld]_{4}\ar[rd]^{4}\\
E&&K(C_2).
}
$$
\end{itemize}

\noindent
In conclusion, we manage to characterize
these families.

\begin{example}
\label{example-chain-of-maps}
In case 1, let $a=b=0$. We get
\begin{equation}
\label{equation-curve-pencil}
C: y^4+xz^3+\underbrace{c}_{\in K\setminus
K^2}x^4=0
\end{equation}
$$
\underbrace{C_1}_{\text{quasiell.}}:xy^2+z^3+
cx^3=0.
$$
Further
$$
\xymatrix{
C\ar[r]&C_1\ar[r]&C_2 \cong\mathbb{P}^1\ar[r]
&C_3 \cong \mathbb{P}^1\ar[r]&C_4\cong
\mathbb{P}^1\ar[r]&\cdots.
}
$$
\end{example}

\medskip\noindent
With $C$ one can construct a pencil
(fibration over $\mathbb{P}^1$)
of quartics.
Let $k=\overline{k}$ (new base field from
what we fixed at the beginning!).
Idea: replace $c$ for $t$ Equation
\ref{equation-curve-pencil}.
 $K=k(\mathbb{P}^1)$,
 $S=V(t_0(y^4+xz^3)+t_1x^4)
\subseteq \mathbb{P}^2_{(x:y:z)}\times
\mathbb{P}_{(t_0:t_1)}$.
This has a singular point, so we blow up:
$$
\xymatrix{
\tilde{S}\ar[d]\ar@/^1em/[dd]^f\\
S\ar[d]\\\mathbb{P}^1
}
$$
The generic fiber is the curve
$y^4+xz^3+tx^4=0$,
so that $t_1/t_0 \in k(\mathbb{P}^1)=K$.
Fibers are plane rational quartics,
and the singular fiber $f_1(0:1)$ is a
configuration of lines
given by a Dynkin diagram
$$
\xymatrix{
E_1 \ar@{-}[r]& E_2\ar@{-}[r]
& E_3 \ar@{-}[r]& E_4\ar@{-}[d]\ar@{-}[r]
& \cdots \ar@{-}[r]& E_{15}\\
& & & E
}
$$
We do the same with $C_1$: define
$S^1=V(t_0(xy^2+z^3)+tx^3) \subseteq
\mathbb{P}^2 \times \mathbb{P}^1$.
Then $S_1$ is also singular, and we blow up
to obtain a smooth
quasielliptic fibration ($\tilde{S^1}$ is
quasielliptic).
The fibers are plane cuspidal cubics.
The singular fiber an arrangement of curves
$$
\xymatrix{
F_1 \ar@{-}[r]& F_2\ar@{-}[r]&
F_3\ar@{-}[d]\ar@{-}[r]& \cdots \ar@{-}[r]&
F_8\\
& & F
}
$$

In fact, there is a correspondence between
the
singular fibers of the two fibrations.

\section{Variations on a theorem of Baum and
Bott}
\label{section-baum-and-bott}

\noindent
Stéphane Druel, Université Claude Bernard
Lyon.
Algebra Seminar, IMPA.
November 5, 2025.

\medskip\noindent
{\bf Abstract.}
In this talk, I will discuss an analogue of
the Baum-Bott vansihing theorem for
the tangent bundle of the space of leaves of
regular foliations together with
applications to transversely projectively
flat foliations of dimension one. This
is work in progress.

\medskip\noindent



Baum-Bott 70'.

Let $E$ be a vector bundle on $X$
complex manifold is a sheaf of 1-jets of $E$:
$$
J'_X(E)= E \oplus \Omega^1_X \otimes E
$$
where equality is the category of sheavs of
abelian groups.
Where the $\mathcal{O}_X$ structure is given
by
$$
\underbrace{f}_{\substack{\text{local
function} \\ \text{on }U}}
(\underbrace{e}_{\substack{\text{local
section} \\ \text{of $E$ on  $U$}}}
, \underbrace{\alpha}_{\substack{\text{local
section} \\
\text{of $\Omega^1_X \otimes E$}}})
=(fe,f\alpha-df \otimes e)
$$
The {\it Atiyah class} of $E$ is the element
$$
\text{at}_X(E) \in H^1(X,\Omega_X^1 \otimes
\text{End}(E))
$$
corresponding to the exact sequence of
$\mathcal{O}_X$-modules
$$
\xymatrix{
0\ar[r]&\Omega_X^1\otimes
E\ar[r]&J_X^1(E)\ar[r]&E\ar[r]&0.
}
$$

\begin{remark}
\label{remark-holomorphic-connection}
 $\text{at}(E)=0$ if and only if there exists
$\nabla:E \to \Omega^1_X \otimes E$ map of
abelian sheaves
such that $e \mapsto  (e-\nabla e)$ is
$\mathcal{O}_X$-linear
$$
-\nabla fe=-f\nabla e - df \otimes e
$$
if and only if $E$ admits a holomorphic
connextion.
\end{remark}

\begin{example}
\label{example-line-bundl-holom-conne}
If $E$ is a line bundle, at
$(E)=(d\text{log}f_{ij})_{ij}
\in H^1(X,\Omega_X^1)$, $E=(f_{ij})$.
\end{example}

\begin{theorem}[Atiyah, '57]
\label{theorem-atiyah-57}
Let $r =\text{rank}E$, $k$ an integer, $X$ be
compact Kähler
and $\Sigma_k$ a polynomial of degree $k$ on
$M_r$
such that  $\Sigma_k(g)$ is the $k$-th
elementary
symmetric function on the characteristic
zeros of
$g \in M_r$. $S_k=$ polynomial given by
polarization.
$$
\begin{aligned}
M_r^{\otimes k}&\to \mathbb{C}\\
\text{End}E^{\otimes k}& \to \mathcal{O}_X\\
\text{Ex., $k=1$}\qquad
\text{End}E&\xrightarrow{\text{Tr}}
\mathcal{O}_X.
\end{aligned}
$$

$$
\xymatrix{
H^1(X,\Omega_X^1 \otimes
\text{End}(E))^{\otimes k}\ar[r]^{\cup }
&H^k((\Omega^1_X)^{\otimes k}\otimes
\text{End}E^{\otimes k}\ar[r]
&H^k(X,\Omega_X^k)\ar[r]
&H^{2k}(X,\mathbb{C})\\
\text{at}(E)^{\otimes k}\ar[rrr]
&&&(2\pi i)^kc_{2k}(E).
}
$$
(Double check the term $(\Omega^1_X)^{\otimes
k}$…
I think there's a typo.)

Partial connection. $\mathcal{F}$ regular
distribution.
$\mathcal{F} \subseteq T_X$ subbundle.
An {\it $\mathcal{F}$-connection} on $E$ is
$\nabla:E \to \mathcal{F}^\vee \otimes E$
such that
$$
\nabla(fe)=f\nabla e+f_{\mathcal{F}}f \otimes
e
$$
where
$d_{\mathcal{F}}:\mathcal{O}_X\xrightarrow{d}
\Omega_X^1
\to\mathcal{F}^\vee$.
\end{theorem}

\begin{proposition}[Bott-Baum, '70]
\label{proposition-bott-baum}
If $E$ admits an $\mathcal{F}$-connection
then all characteristic
classes of degree  $>\text{codim}\mathcal{F}$
vanish.
\end{proposition}

\begin{proof}
There exists an $\mathcal{F}$-connection if
and only if
$$
\xymatrix{
0\ar[r]&\Omega_X^1 \otimes
X\ar[r]\ar[d]&J_X^1\ar[r]&E\ar@{=}[d]\ar[r]&
0\\
0\ar[r]&\mathcal{F}^\vee\otimes E\ar[r]&
\mathcal{F}^\vee \otimes E \oplus E\ar[r]&
E\ar[r]&0
}
$$
if and only if $\text{at}(E) \in
H^1(X,\Omega_X^1\otimes\text{End}E)$
maps to 0 in
$H^1(X,\mathcal{F}^\vee\otimes\text{End}E)$
is and only if $\text{at}(E)$ comes from
$$
H^1(X,N^\vee\otimes\text{End}E)\to
H^1(X,\Omega_X^1\otimes\text{End}E)
$$
where $N=T_{X/\mathcal{F}}$.
\end{proof}

\begin{example}[Bott's partial connection]
\label{example-bott-partial-connection}
Suppose $\mathcal{F}$ is a regular foliation
and $N=T_{X/\mathcal{F}}$ has a
$\mathcal{F}$-connection.
Take $n$ a local section of $N$,  $v$ a local
section of $\mathcal{F}$
(over the same open) and suppose exists a
section $t$ of $T_X$
that is a lift of $t$, i.e. $p(t)=n$ for
$p:T_X \to N$.
Then $\nabla_v^Bn=p([v,t])$ is a section of
$N$.
This defines an $\mathcal{F}$-connection on
$N$.
\end{example}

\begin{lemma}[Beauville]
\label{lemma-beauville}
Let $X$ be a complex manifold and $E$ is a
direct submanifold of $T_X$.
Then $\text{at}(E) \in H^1(X,\Omega_X^1
\otimes \text{End}E)$
comes from $H^1(X,E^\vee \otimes
\text{End}(E))$.
\end{lemma}

\begin{proof}
Consider
$$
T_X= E \oplus F \xrightarrow{p}E.
$$
We claim that there exists an
$\mathcal{F}$-connection on $E$.
Let $v$ be a local section of $F$,
$e$ a local section of $E$.
Define $\nabla_ve=p([v,e])$.
One checks that this satisfies the required
properties.
\end{proof}

\noindent
Notice that no involutivity is imposed on $E$
- it's just a distribution.

\medskip\noindent
Now we will discuss some variations of
Baum-Bott's theorem
\ref{theorem-baum-bott}.

Let $\mathcal{F}$ be  regular foliation, $E$
a vector bundle on $X$ equipped with a flat
$\mathcal{F}$-connection.
For example, $\nabla^B$ flat.

Recall that  $(E,\nabla)$ can be thought of
as
$\text{Mon}(\mathcal{F})$-equivariant vector
bundles
where $\text{Mon}(\mathcal{F})$ is the
monodromy groupoid.

Let $X/\mathcal{F}$ be the space of leaves,
and
$\mathcal{O}_{X/\mathcal{F}}=\mathcal{O}
_X^{d_{\mathcal{F}}}=
\{\text{functions $f$ such that
$d_{\mathcal{F}}f=0$}\}$.

It is a fact that $E^{\nabla}$ is a locally
free $\mathcal{O}_{X/\mathcal{F}}$
sheaf of rank the rank of $E$ and such that
$E=E^{\nabla}\otimes_{F_{X/\mathcal{F}}}
\mathcal{O}_X$.

Now one defines bundles of 1-forms and
tangent bundle as follows.
 $\Omega_{X/\mathcal{F}}^1=(N_{\mathcal{F}}^*
 )^{\nabla B}$
and
$T_{X/\mathcal{F}}=N_{\mathcal{F}}^{\nabla
B}$,
which are locally free sheaves of
$\mathcal{O}_{X/\mathcal{F}}$-modules
of rank $q=\text{codim}\mathcal{F}$.

The Atiyah class of $(E,\nabla)$ is
$$
J^1_{X/\mathcal{F}}(E,\nabla)=E^{\nabla}
\oplus
\Omega_{X/\mathcal{F}}^1
\otimes_{\mathcal{O}_{X/\mathcal{F}}}
E^{\nabla}
$$
with $\mathcal{O}_{X/\mathcal{F}}$-module
structure given by
 $$
(e,\alpha)=(f\alpha,f\alpha-\underbrace{df}
_{\in
\Omega^1_{X/\mathcal{F}}}
\otimes e.
$$
Then
$$
\text{at}_{X/\mathcal{F}}(E,\nabla)
\in
H^1(X,\Omega_{X/\mathcal{F}}
^1\otimes_{\mathcal{O}_{X/\mathcal{F}}}
\text{End}_{\mathcal{O}_{X/\mathcal{F}}}
E^{\nabla})
$$
corresponding to
$$
\xymatrix{
0\ar[r]&\Omega_{X/\mathcal{F}}
^1\otimes_{\mathcal{O}_{X/\mathcal{F}}}\ar[r]
&J_{X/\mathcal{F}}^1(E)\ar[r]&E^{\nabla}
\ar[r]&0.
}
$$

\begin{lemma}
\label{lemma-atiya-class-maps-atiya-class}
$\text{at}_{X/\mathcal{F}}(E,\nabla)$ maps to
$\text{at}_X(E)$ under
$$
H^1(X,\Omega_{X/\mathcal{F}}
^1\otimes_{\mathcal{O}_{X/\mathcal{F}}}
\text{End}_{\mathcal{O}_{X/\mathcal{F}}}
E^{\nabla}
\to H^1(X,\Omega_X \otimes \text{End}E).
$$
\end{lemma}

\noindent
Partial connections. Let $G \subseteq
T_{X/\mathcal{F}}$
be a subbundle a {\it $G$-connection}
$(E,\nabla)$ is
$$
\begin{aligned}
D: E^{\nabla} &\longrightarrow G^\vee\otimes
E^{\nabla} \\
\text{s.t.} D(fe)&=fDe+d_Gf\otimes e
\end{aligned}
$$
where
$d_G:\mathcal{O}_{X/\mathcal{F}}\to\Omega_{X/
\mathcal{F}}^1\to
G^\vee$.

It's a fact that a $G$-connection
$(E,\nabla)$ exists if and only if
$\text{at}_{X/\mathcal{F}}(E,\nabla)$ maps to
0 in
$H^1(X,G^\vee
\otimes_{G_{X/\mathcal{F}}}\text{End}
_{\mathcal{O}_{X/\mathcal{F}}}
E^{\nabla}$.

The analogue of Beauville's lemma
\ref{lemma-beauville} is

\begin{lemma}[Beauville]
\label{lemma-analogue-beauville}
Suppose $N_{\mathcal{F}}=\mathcal{A} \oplus
B$ such that
$\nabla^B\mathcal{A} \subseteq
\mathcal{F}^\vee
\otimes_{\mathcal{O}_X}\mathcal{A}$
and $\nabla^B \mathcal{B} \subseteq
\mathcal{F}^\vee\otimes_{\mathcal{O}_X}
\mathcal{B}$.
Let $\nabla_{\mathcal{A}}$ and
$\nabla_{\mathcal{B}}$ be a flat
 $\mathcal{F}$-connection on $\mathcal{A}$
 and $\mathcal{B}$.

Then
$\text{at}_{X/\mathcal{F}}(\mathcal{A},
\nabla_{\mathcal{A}})
\in
H^1(X,\Omega_{X/\mathcal{F}}
^1\otimes_{\mathcal{O}_{X/\mathcal{F}}}
\text{End}_{\mathcal{O}_{X/\mathcal{F}}}
\mathcal{A}^{\nabla_{\mathcal{A}}}$
comes from
$H^1(X,(\mathcal{A}^{\nabla_{\mathcal{A}}})
^\vee
\otimes_{X/\mathcal{F}}\text{End}
_{\mathcal{O}_{X/\mathcal{F}}}\mathcal{A}
^{\nabla_{\mathcal{A}}}$.
\end{lemma}

\begin{proof}
Same proof but we need a Lie bracket on
$T_{X/\mathcal{F}}=N_{\mathcal{F}}\subseteq
N_{\mathcal{F}}\overset{p}{\leftarrow}T_X$.
Let $v_1,v_2$ be local sections on
$T_{X/\mathcal{F}}$
with $v_1=p(t_1)$ and $v_2=p(t_2)$.
Define
$$
[v_1,v_2]=p([t_1,t_2]).
$$
One can check that this a Lie bracket on
$T_{X/\mathcal{F}}$.
Suppose that $t_1 \in \mathcal{F}$. Then
$$
p([t_1,t_2])\overset{\text{def}}{=}
\nabla_{t_1}^{\mathcal{B}}v_2=0.
$$
For $v_3 \in \mathcal{F}$,
$$
\nabla_{v_3}^{\mathcal{B}}[v_1,v_2]
=p([v_3,[t_1,t_2]])=\pm\nabla_{v_3}
^{\mathcal{B}}v_1\pm
\nabla_{v_3}^{\mathcal{B}}v_2=0.
$$
\end{proof}

\begin{lemma}
\label{lemma-vanishing-chern-class}
Suppose that $\mathcal{F}$ is a regular
foliation on $X$.
Let $N_{\mathcal{F}}=\mathcal{L}_1 \oplus
\ldots \oplus \mathcal{L}_q$
be line bundles on $X$,
$\nabla^{\mathcal{B}}\mathcal{L}_i\subseteq
\mathcal{F}^\vee\otimes\mathcal{L}
_i\rightsquigarrow
\nabla_i$.
Suppose
$(\mathcal{L}_i,\nabla_i)\cong(\mathcal{L}_j,
\nabla_j)$
for all $i,j$.
Then $c_1(\mathcal{L}_i)=0$.
\end{lemma}

\begin{proof}
By the Comparison Lemma,
$$
\begin{aligned}
H^1(X,\Omega_{X/\mathcal{F}}^1
&\longrightarrow H^1(X,\Omega_X^1) \\
\text{at}_{X/\mathcal{F}}(\mathcal{L},
\nabla_i)
&\longmapsto
\text{at}(\mathcal{L}_i),
\end{aligned}
$$

\noindent
it suffices to prove that
$\text{at}_{X/\mathcal{F}}(\mathcal{L}_i,
\nabla_i)$
vanishes.
$$
\text{at}_{X/\mathcal{F}}(\mathcal{L}_i,
\nabla_i)
\in H^1(X,\Omega_{X/\mathcal{F}}^1)
=\bigoplus_{j=1}^9
H^1(X,(\mathcal{L}_j^{\nabla_j})^\vee)
$$
comes from
$H^1(X,(\mathcal{L}_i^{\nabla_i})^\vee)$.
\end{proof}

\begin{example}
\label{example-degree-2-line-bundle-on-e}
$A=E^1 \times B^{n-1}$, an abelian variety.
Let $ L$ be the pullback of a degree 2
divisor on $E$.
Then there exists  $T_X
\to\!\!\!\!\!\to L^{\oplus n-1}$.
\end{example}

\begin{proposition}
\label{proposition-}
$\mathcal{F}$ regular foliation of dimension
1 on $X$ complex projective
of dimension $\geq 3$.
Suppose $\mathcal{F}$ is transversally
projective flat
and that $X$ contains a rational curve.
Then $\mathcal{F}$ is a $\mathbb{P}^1$-bundle
and it is tangent
to $\mathcal{F}$.
\end{proposition}

\medskip\noindent
Now we explain what it means to be
projectively flat.
It is when the projectivization of the normal
bundle $\mathbb{P}(N)$
is induced by a representation of $\pi_1(X)
\to \text{PGL}_g$
it comes with an Ehresmann connection on
$$
\xymatrix{
\mathbb{P}(N)=\mathbb{P}\ar[d]^{f}\\
X
}
$$
$T_{\mathbb{P}}=T_f \oplus g$.

Now, it is also called transversally
projectively flat
if $\nabla^B$ is a partial connection on $N$
and there exists a foliation $\mathcal{H}$ on
$$
\xymatrix{
\mathbb{P}(N)\ar[d]\\ X
}
$$
such that $\mathcal{H}$ projects onto
$\mathcal{F}$.
Transversality means that $\mathcal{H}
\subseteq \mathcal{G}$.

Finally, transversally projectively flat
means that
it is induced on a
$\text{Mon}(\mathcal{F})$-module
by a representation of
$\pi_1(\text{Mon}(\mathcal{F}) \to
\text{PGL}_g$.

\section{Um passeio nos trabalhos de
Gerd Faltings}
\label{section-faltings}

\noindent
Hossein Movasati, IMPA.
Pelestra especial,
May 19, 2026.

\medskip\noindent
{\bf Abstract.}
Nesta palestra vamos dar um passeio nos
trabalhos de G. Faltings, começando por sua
demonstração da conjetura de Mordell em
1983, que lhe trouxe a medalha Fields. Uma
boa parte de sua carreira foi elaborar o
lado aritmético de um dicionário heurístico
proposto por P. Deligne em 1970, que tem
origem na teoria conjetural de motivos de
Grothendieck. Isso inclui versão p-ádica de
teoria de Hodge, domínios de períodos,
correspondência de Simpson, etc. Várias
propriedades aritméticas de variedades
abelianas são outro foco da carreira dele
que vamos abordar nesta palestra. Em 2026
ele ganhou o prêmio Abel por suas
contribuições fundamentais em geometria
aritmética.

\medskip\noindent
{\bf Upshot.} 
Elliptic curves are diophantine equations.
Some are simple and fun to solve,
some are extremely complicated ---
like Fermat's last theorem.
Finding solutions to a diophantine
equation is like finding $k$-points
on an algebraic curve for suitable
choices of fields.
In this sense,
Mordell's conjecture, now Faltings'
theorem is a statement of finiteness
of solutions of diophantine equations.
Other results are adaptations
of results from complex geometry
to finite fields.

\medskip\noindent
{\bf Notes.}

Finding roots of a polynomial $f \in
\mathbb{Q}[x,y]$ becomes the geometric
problem of finding the $\mathbb{Q}$-rational
points of $\operatorname{Spec}(\mathbb{Q}[x,y]/(f))$.
Or, one can pass to the projective version
$\operatorname{Proj}(\mathbb{Q}[x,y,z]/(F))$,
where $F$ is the homogenization of $f$.
We write
$C(\mathbb{Q})=\{[x:y:z] \in
\mathbb{P}^2(\mathbb{Q}) \mid
F(x,y,z)=0\}$
and similarly $C(\mathbb{C})$ for the
analogous set.
 
It turns out that the topology
of the curve, via its genus, strongly
constrains the $\mathbb{Q}$-rational points.
The easiest cases are:
\begin{itemize}
\item ($g=0$.) Here $d=1,2$;
then the curve is projectively equivalent to
a conic, and over $\mathbb{Q}$ it has either
no rational points or infinitely many.

\item ($g=1$.) Here $d=3$; this is more difficult
but still well-studied; see the
Birch--Swinnerton-Dyer conjecture. 

\item ($g\geq 2$.) Here $d\geq 4$.
This is much more complicated!
Faltings'
theorem implies that $C(\mathbb{Q})$ is
finite.
\end{itemize}

\noindent
This is the former Mordell conjecture
from 1962.

\begin{theorem}[Faltings, 1983]
For any smooth curve of genus $\geq 2$ over
$\mathbb{Q}$, we have $\#C(\mathbb{Q})<
\infty$.
\end{theorem}

The same theorem remains true after
replacing $\mathbb{Q}$ by a number field,
that is, a finite field extension of
$\mathbb{Q}$.

\begin{theorem}[Manin 1963, Grauer 1965]
\label{theorem-manin-grauer}
For any smooth non isotrivial curve of
genus $\geq 2$ over $k=\mathbb{C}(t)$
we have
$$
\# C(k) < \infty.
$$
\end{theorem}

\begin{proof}
Using Gauss-Manin connection ---
the name of which was invented by 
Grothendieck.
\end{proof}

\noindent
The main difference between
$\mathbb{C}(t)$
and $\mathbb{Q}$
is that in $\mathbb{C}(t)$
we have a derivation
$$
\frac{\partial }{\partial t}:
\mathbb{C}(t) \to \mathbb{C}(t)
$$
satisfying the Leibniz identity.

\medskip\noindent
The proof is somehow related
to Shafarevich and Tate conjectures.
The former states there's only
finitely many isomorphism
classes of abelian varieties of a
given dimension with
fixed polarization degree, defined
over a field with
good reduction outside a finite set
of places.

Tate conjecture was proved for finite
fields by Tate in 1966,
and by Faltings for number fields
in 1983.
It relates maps between the tangent
spaces of two abelian varieties
and maps between the varieties
themselves:
$$
\Hom_K(A,B) \otimes \mathbb{Z}_p
\cong
\Hom_{Gal(\tilde{K}/K)}
(T_pA,T_pB)
$$
where $K$ is a field of characteristic
$\neq p$, and $A,B$ are abelian over $K$.

Further, the proof uses Arakelov's
theory of arithmetic genus
and Faltings height.

\medskip\noindent
{\bf Abelian varieties.} 

Any curve (over any field)
can be embedded in some variety
called the Jacobian variety.
Over $\mathbb{C}$,
abelian varieties are just tori:
$$
\mathbb{C}^n/\Lambda
\cong_{C^\infty}
S^1\times… S^1,\qquad  2n \text{ times}.
$$

An abelian variety of dimension one
over a field $k$ is called
an {\it elliptic curve}.
They are given by diophantine
equations of degree 3 in two variables:
$$
y^2+a_1xy+a_3y
=
x^3+a_2x^2+a_4x,
\qquad a_1,…,a_6 \in k.
$$

The name was invented by A. Weil.
Before him, many mathematicians like E.
Picard,
G. Humbert, were using ``hyperelliptic
surface'' for Abelian surfaces.

\medskip\noindent
{\bf Shavarevich's conjecture.} 
What is the number of non-isomorphic elliptic
curves over $\mathbb{Q}$ and with bad
reduction only for primes $2,3$?

\medskip\noindent
A consequence of Tate conjecture
and Chebotarev density theorem is:;

\begin{theorem}
\label{theorem-tate-chebot-consequence}
LEt $E_i$, $i=1,2$ be two
elliptic curves over $\mathbb{Q}$.
TFAE:
\begin{enumerate}
\item $E_1$ is isogeneous to $E_2$
over $\mathbb{Q}$.
\item For all except a finite number
of primes,
$$
\# E_1(\mathbb{F}_p)
=\# E_2(\mathbb{F}_p).
$$
\end{enumerate}
\end{theorem}

\begin{definition}
\label{definition-isogeny}
{\it AI-generated.}
An isogeny between elliptic curves over
$\mathbb{Q}$ is a nonzero morphism of
elliptic curves defined over $\mathbb{Q}$
that is also a group homomorphism and has
finite kernel.
\end{definition}

Here isogenous means that there exists an
isogeny $f:E_1 \to E_2$ over $\mathbb{Q}$.

\medskip\noindent
One of the key ingredients
in Falting's proof
is the notion of {\it Faltings height},
which are integrals used to
control the numerator and denominator
or rational solutions of a (diophantine?)
equation.

For example, for $r = n/m$,
$h(r)=n^2+m^2$. For all $M \in \mathbb{R}$,
$$
\{r \in \mathbb{Q}| h(r)<M\}
$$
is finite.

\begin{proposition}[Main property
of height functions]
\label{proposition-faltings-height}
For any positive integer $M$,
the number of elliptic curves $E/\mathbb{Q}$
with $h_{\text{Falt}}<M$ is finite.
\end{proposition}

\noindent
On the curve given by Equation
\ref{equation-?}
consider the 1-form
$$
\omega=\frac{dx}{2y+a_1x+a_3}.
$$
The Falting's height is basically…
[missing! an integral involving $\omega$…]

\medskip\noindent
\begin{itemize}
\item Elliptic integrals $\leq 1900$.
\item Poincaré, Pircard, $\sim 1900$.
\item Singular, de Rham 1900-1950.
\item Weil: cohomology over the integers
(enter arithmetic).
\item Grothendieck: étale, crystalline, etc.
\end{itemize}

\noindent
Then is Deligne's dictionary,
in a 1970 paper, where $\ell$-adic
cohomology is introduced
(I think --- abstract in french, check).

\medskip\noindent
From complex to arithmetic geometry:

\begin{itemize}
\item Endilchkeitssätze
fur abelsche Varietäten über Zahlkörpern.
\item 1984: Calculus on arithmetic surfaces.
\item 1988: p-adic Hodge theory.
\item 1990: Degeneration of abelian
varieties (with Chai Ching-li).
\item 1992: Lectures on the
arithmetic Riemann-Roch theorem.
\item 1999: Does there exist an arithmetic
Kodaira-Spencer class? (!!!)
\item 2005: A p-adic Simpson corresponding.
\item 2020: covering of p-adic period
domains.
\end{itemize}

\medskip\noindent
{\bf p-adic Hodge theory.}

Grothendieck and Deligne
noticed that we can obtain
the Hodge decomposition of cohomology
using filtrations, which sit
in the realm of polynomials ---
no harmonic analysis needed.

The idea of Deligne (no, Faltings?) 
was to use
this proof for the $p$-adic case.

\section{Grothendieck--Teichm\"uller theory: from arithmetic geometry to arithmetic topology}
\label{section-grothendieck-teichmuller-arithmetic-topology}

\noindent
Marco Boggi.
Qua 27 mai 2026, 15:30.
SALA 228.

\medskip\noindent
{\bf Abstract.}
Grothendieck--Teichm\"uller theory originated
in arithmetic geometry through the study of
Galois actions on geometric \'etale
fundamental groups. In this talk, we argue
that arithmetic topology provides a natural
setting for applications of
Grothendieck--Teichm\"uller theory. More
specifically, we explain how our theory of
automorphisms of profinite mapping class
groups may lead to new results on profinite
rigidity questions for mapping class groups.

\medskip\noindent
{\bf Warning.}
These are live notes from the talk,
written with unusually heavy AI assistance
since that day I forgot my glasses.
Some theorem
statements may be imprecise; the goal is to
record the ideas and references from the day.\medskip\noindent

\medskip\noindent
The main upshot is that GT theory allows us
to pass the problem to a free group in three
generators, and that the talk aims to
contrast arithmetic and topological
properties.

A more personally significant 
upshot, explained next,
is that the étale fundamental
group of the moduli space of curves
is the profinite completion of
the fundamental group of any curve in 
the moduli, which makes sense since the
underlying topological space of all curves
in the moduli is the same.

\medskip\noindent
Historically, this is very much in the spirit
of Grothendieck's 1984 \emph{Esquisse d'un
programme}. Drinfeld then introduced the
Grothendieck--Teichm\"uller group in 1991,
making precise a framework that was already
expected to have arithmetic applications.

\medskip\noindent
The point is that $\mathcal M_{g,n}$ is
not just an abstract moduli space: it is
built from one fixed topological surface
$S_{g,n}$. A point of $\mathcal M_{g,n}$
is essentially a complex structure on this
same underlying surface.

\medskip\noindent
Thus the topological fundamental group of
the moduli space records how a family of
curves can move around and come back to the
same curve with its underlying surface
changed by a diffeomorphism. This gives the
mapping class group:
$$
\pi_1^{\mathrm{top}}
\bigl(\mathcal M_{g,n}(\mathbb C)\bigr)
\cong
\Gamma_{g,n}.
$$
Here
$$
\Gamma_{g,n}
=
\operatorname{Diff}^+(S_{g,n})/
\operatorname{Diff}_0(S_{g,n}).
$$

\medskip\noindent
The \'etale fundamental group is the
algebraic-geometric version of the
fundamental group. Instead of remembering
all covering spaces, it remembers finite
algebraic covers. Over $\mathbb C$, this
amounts to taking the profinite completion
of the usual fundamental group. Therefore
one gets
$$
\pi_1^{\acute{e}t}
\bigl(\mathcal M_{g,n}
\otimes_{\mathbb Q}\overline{\mathbb Q}\bigr)
\cong
\widehat{\Gamma}_{g,n}.
$$

\medskip\noindent
So the profinite mapping class group is not
an artificial object: it is literally the
\'etale fundamental group of the moduli
space of curves. This is why
Grothendieck--Teichm\"uller theory naturally
acts on profinite mapping class groups.

\medskip\noindent
The starting point is the \'etale
fundamental group of the projective line with
three points removed. Over $\mathbb{C}$, the
space
$$
\mathbb{P}^1_{\mathbb{C}} \setminus \{0,1,\infty\}
$$
is a sphere with three punctures, so its
topological fundamental group is a free group
on two generators:
$$
\pi_1^{\mathrm{top}}
\bigl(\mathbb{P}^1_{\mathbb{C}} \setminus
\{0,1,\infty\}\bigr)
\cong F_2.
$$
After profinite completion one gets
$$
\pi_1^{\acute{e}t}
\bigl(\mathbb{P}^1_{\overline{\mathbb{Q}}} \setminus
\{0,1,\infty\}\bigr)
\cong \widehat{F}_2,
$$
the free profinite group on two generators.

\medskip\noindent
Now let
$$
X=\mathbb{P}^1_{\mathbb{Q}} \setminus
\{0,1,\infty\}.
$$
There is an exact sequence of profinite
groups
$$
1\to
\pi_1^{\acute{e}t}(X_{\overline{\mathbb{Q}}})
\to
\pi_1^{\acute{e}t}(X)
\to
\operatorname{Gal}(\overline{\mathbb{Q}}/\mathbb{Q})
\to 1.
$$
Because the geometric fundamental group is
\(\widehat{F}_2\), the absolute Galois group
acts by outer automorphisms on a free
profinite group:
$$
\operatorname{Gal}(\overline{\mathbb{Q}}/\mathbb{Q})
\longrightarrow
\operatorname{Out}(\widehat{F}_2).
$$
This is one of the first places where
Grothendieck's vision becomes concrete.
The theorem says that this action is
faithful.

\medskip\noindent
In particular, the absolute Galois group can
be studied through its action on the
combinatorics of finite coverings of
\(\mathbb{P}^1\setminus\{0,1,\infty\}\).
This is the origin of dessins d'enfants.
A finite cover of
\(\mathbb{P}^1\setminus\{0,1,\infty\}\)
corresponds to a finite quotient of
\(\widehat{F}_2\), and the Galois action
permutes these quotients in a very rigid way.
That rigidity is what eventually leads to the
Grothendieck--Teichm\"uller group.

\medskip\noindent
The point of the talk is that a similar
picture appears in arithmetic topology.
Mapping class groups also act on profinite
fundamental groups of surfaces, and one can
study their automorphisms in a way that
mirrors the Galois action above. This
suggests profinite rigidity questions: to
what extent does the profinite completion of
a mapping class group remember the original
group? The Grothendieck--Teichm\"uller
framework provides a natural language for
these problems, and may lead to new rigidity
results.

\medskip\noindent
In the arithmetic case there are more
automorphisms than in the topological case,
and this is what makes the theory useful for
applications. The speaker suggested that
these methods should lead to a proof of a
failure of profinite rigidity: more
precisely, if \(d(S)\) denotes the complexity
of the surface, then there should exist
finite-index subgroups \(\Gamma\) and
\(\Gamma'\) of \(\Gamma(S)\) such that
$$
\widehat{\Gamma}\cong \widehat{\Gamma'}
\qquad\text{but}\qquad
\Gamma\not\cong \Gamma'.
$$

\medskip\noindent
For \(g=0\), the speaker can prove that every
pointwise inner automorphism of
\(\widehat{\Gamma}(S)\) is actually inner.
That is, if an automorphism acts on every
single element as conjugation, then the
automorphism itself is conjugation.

\medskip\noindent
As a corollary, the Grothendieck--Teichm\"uller
group acts faithfully on the conjugacy
classes of elements of the profinite
completion of the braid group
\(\widehat{B}_n\). This is interesting from
the topological point of view because of the
Alexander theorem, which characterizes
braids via their relationship with links.

\medskip\noindent
Another application mentioned in the talk is
Furusho's work on finite-type (Vassiliev)
knot theory. The GT action can be used to
organize knot-theoretic information in a way
that is compatible with the profinite braid
group viewpoint.

\medskip\noindent
Let $B$ denote the profinite mapping class
group under consideration. We write
\(\operatorname{Aut}^{\mathbb{I}_0}(B)\) for
the subgroup of \(\operatorname{Aut}(B)\)
consisting of those automorphisms that
preserve the procyclic subgroups generated by
the elements \(C_r\) arising from Dehn
twists. The centralizer of the natural image
of the Grothendieck--Teichm\"uller action in
\(\operatorname{Aut}^{\mathbb{I}_0}(B)\) is
what we call the Grothendieck--Teichm\"uller
group \(G(c,0)\).

\medskip\noindent
For a connected oriented surface
\(S_{g,n}\) of genus \(g\) with \(n\)
punctures and \(\chi(S_{g,n})<0\), the
mapping class group is
the group of orientation-preserving
diffeomorphisms modulo isotopy to the
identity. More precisely, if
\(\Gamma_{g,n}\) denotes the mapping class
group of \(S_{g,n}\), then
\(\Gamma_{g,n}=\mathrm{P}\Gamma(S_{g,n})\)
in the pure case.

\begin{theorem}\label{thm:mapping-class-moduli-space}
There is a natural isomorphism
$$
\Gamma_{g,n}\cong
\pi_1^{\mathrm{top}}
\bigl(\mathcal{M}_{g,n}(\mathbb{C})\bigr),
$$
where \(\mathcal{M}_{g,n}\) is the moduli
space of smooth genus \(g\) curves with
\(n\) marked points.
\end{theorem}

\begin{proof}
This is the classical identification of the
mapping class group with the orbifold
fundamental group of the moduli space. In the
case of ordered marked points, the pure
mapping class group corresponds to the
fundamental group of the pointed moduli
space.
\end{proof}

\medskip\noindent
There is also a standard exact sequence
relating the pure mapping class group to the
full mapping class group:
$$
1\to \widehat{\mathrm{P}\Gamma}_{g,n}
\to \widehat{\Gamma}_{g,n}
\to \Sigma_n
\to 1.
$$
At the discrete level this comes from the
action of the mapping class group on the set
of punctures, and the quotient records the
induced permutation of the punctures.

\begin{remark}
How can we study the Deligne--Mumford
boundary through points on a surface? This
appears to be a standard method discussed in
the talk.
\end{remark}

\medskip\noindent
The theorem implies that the \'etale
fundamental group of the moduli space is the
profinite completion of the mapping class
group:
$$
\pi_1^{\acute{e}t}
\bigl(\mathcal{M}_{g,n}\otimes_{\mathbb{Q}}\overline{\mathbb{Q}}\bigr)
\cong
\widehat{\Gamma}_{g,n}.
$$
In the pure case this is
\(\widehat{\mathrm{P}\Gamma}(S_{g,n})\).

\begin{theorem}[Hoshi--Mochizuki]\label{thm:hoshi-mochizuki-centralizer}
Let \(\Sigma_n\) act on the profinite
mapping class group in the natural way by
permuting the punctures. Then there is a
natural injection
$$
Z_{\operatorname{Out}^{\mathbb{I}_0}
(\widehat{\Gamma}_{g,n})}(\Sigma_n)
\hookrightarrow
Z_{\operatorname{Out}^{\mathbb{I}_0}
(\widehat{\Gamma}_{g,n+1})}(\Sigma_{n+1}),
$$
and this map is an isomorphism for \(n\ge 5\).
\end{theorem}

\begin{proof}
This is the centralizer comparison used in
the study of profinite mapping class groups
and their puncture-permutation symmetries.
The low-puncture cases require separate
analysis, while for \(n\ge 5\) the natural
stabilization map becomes an isomorphism.
\end{proof}

\begin{theorem}[Hoshi--Minamide--Mochizuki]\label{thm:hoshi-minamide-mochizuki}
There is a natural isomorphism
$$
\operatorname{Out}^{\mathbb{I}_0}
\bigl(\mathrm{P}\check\Gamma(S)\bigr)
\cong
\Sigma_n \times \widehat{\operatorname{GT}}.
$$
Equivalently, in the genus zero punctured
sphere case one may write
$$
\operatorname{Out}^{\mathbb{I}_0}
\bigl(\mathrm{P}\check\Gamma_{0,n}\bigr)
\cong
\Sigma_n \times \widehat{\operatorname{GT}}.
$$
\end{theorem}

\begin{proof}
This is the theorem that removes the extra
rigidity hypotheses and identifies the full
outer automorphism group, under the
\(\mathbb{I}_0\)-condition, as the product of
the puncture permutation group and the
profinite Grothendieck--Teichm\"uller group.
\end{proof}

\medskip\noindent
There is also a slogan in the genus-zero
setting:
$$
\operatorname{Out}\bigl(\widehat{\Gamma}_{0,n}\bigr)
\cong \widehat{\operatorname{GT}}
$$
for the appropriate completion and rigidity
condition encoded by the exponent on
\(\operatorname{Out}\). More generally, for
small genus \(g\leq 2\), the speaker proved a
corresponding isomorphism
$$
\operatorname{Out}^{\ast}\bigl(\widehat{\Gamma}_{g,n}\bigr)
\cong
\Sigma_n \times \widehat{\operatorname{GT}},
$$
where the superscript \(\ast\) indicates the
same rigidity condition as in the previous
theorems.

\section{Hodge structures in toric ambient}
\label{section-hodge-structures-toric-ambient}

\noindent
Roberto Villaflor, USM, Chile.
GADEPS Focused Meeting June 29, 2026, IMPA.

\medskip\noindent
Joint with L. Broring, G. Martinez,
A. Viergever.

\medskip\noindent
{\bf Hodge structures over a perfect
field.} Let $X$ be a smooth
projective variety over $\mathbb{C}$.
We have the Hodge-De Rham spectral sequence
$E^{p,q}_1=H^q(X,\Omega_X^p) 
\implies \mathbb{H}^k(X,\Omega_X^\bullet)
\simeq H^k_{dR}(X/\mathbb{C})$.
It degenerates at $E_1$, giving
a non-canonical isomorphism
$$
H^\bullet_{dR}(X/\mathbb{C})
\simeq
\bigoplus_{p+q=k}H^q(X,\Omega_X^p)
$$
We may ask ourselves whether this is also
true for other fields.

\begin{remark}[Deligne-Illusie]
\label{remark-derham-to-hodge-not-degenerate}
In general, for $X/k$ perfect field,
this may not degenerate at $E_1$.
\end{remark}

\noindent
For $X = \{F=0\} \subseteq \mathbb{P}^{n+1}$
smooth degree $d$ hypersurface over
$\mathbb{C}$,
\begin{enumerate}
\item 
$$
H^q(X,\Omega_X^p)
\simeq \begin{cases}
H^q(X,\Omega_X^p)_{\operatorname{prim}}\qquad
&p \neq q \\
H^p(X,\Omega_X^p)_{\operatorname{prim}}\oplus
\mathbb{C}\theta |_X^p\qquad &
\end{cases}
$$
where $\theta$ is the class of a hyperplane.

\item
$H^q(X,\Omega_X^p)_{\operatorname{prim}}=0$
for $p \neq q \neq  n$ 
(Lefschetz hyperplane theorem).

\item Griffith basis theorem for $p+q=n$.
\end{enumerate}

\begin{remark}
\label{remark-lps}
For $X/k$ smooth projective, $k$ perfect
of characteristic $\neq 2$,
Levine-Pepin-Srinivas '24: (1), (2)
and (3) still hold.
\end{remark}

\medskip\noindent
{\bf Quadratic Euler characterstic.} 
Let $X$ be a CW-complex of dimension $n$.
$$
\chi_{\operatorname{top}}(X)
=\sum_{i=1}^n(-1)^i b_i(X),
$$
where
$b_i=\operatorname{rank}H_i(X,\mathbb{Z})$
is the $i$-th Betti number.

\begin{proposition}
\label{proposition-properties-of-euler}
\begin{enumerate}
\item $\chi_{\operatorname{top}}$
is homoropy invariant.
\item $Y \subseteq X$,
$\chi_{\operatorname{top}}(X\setminus Y)
=\chi_{\operatorname{top}}(X)
-\chi_{\operatorname{top}}(Y)+1$.
\item $F \hookrightarrow E \to B$ fiber
bundle, then $\chi_{\operatorname{top}}(E) =
\chi_{\operatorname{top}}(B)
\chi_{\operatorname{top}}(F)$

\item $\hat{\chi_{\operatorname{top}}}(X
\wedge Y) =
\hat{\chi_{\operatorname{top}}}(X)
\hat{\chi_{\operatorname{top}}}(Y)$.
\end{enumerate}
\end{proposition}

\begin{definition}
\label{definition-stably-homotopic}
$X$ is {\it stably homotopic} to $Y$
if
$$
\Sigma^kX \equiv \Sigma^kY
$$
where $\Sigma^k$ is the smash product
of $k$ 1-spheres.
\end{definition}

\noindent
By the reduced Euler characteristic formula,
if $X$ and $Y$ are stably homotopic,
then $\chi_{\operatorname{top}}(X)
=\chi_{\operatorname{top}}(Y)$.

\medskip\noindent
Morel-Boevodky '99: we have
a stable motivic homotopy category.

There exists an analogous construction
of categorical Euler characteristic
$X$ quasi-projective smooth $k$-scheme.
$\chi(X/k) \in GW(k)$
where $GW(k)$ is the Grothendieck-Witten
ring. In characteristic $\neq 2$, 
bilinear symmetric forms = group
generated by the monoid of isometry
classes of quadratic forms with $\oplus$.

The upshot is that instead
of assigning a number, now we put
a quadratic form to our scheme.

\begin{remark}
\label{remark-real-field-euler}
If $k \subset \mathbb{R}$ is a real field,
the rank of $\chi(X/k) =
\chi_{\operatorname{top}}(X(\mathbb{C}))$,
and the signature of $\chi(X/k)$
is
$\chi_{\operatorname{top}}(X(\mathbb{R}))$.
\end{remark}

\begin{theorem}[Levine '20]
\label{theorem-levine20}
$X/k$ smooth projective, $k$ perfect,
char $\neq 2$, then
\begin{enumerate}
\item If $\dim X$ is odd, then
$\chi(X/k) = mH$, $m \in \mathbb{Z}$,
$$
H=\begin{pmatrix}1&0\\ 0&-1\end{pmatrix}
=\langle 1\rangle+ \langle-1\rangle.
$$
\item If $\dim X$ is even, then
$\chi(X/k)=mH+Q$, $m \in \mathbb{Z}$
and $Q$ is determined by the first
page of the Hodge-to-de-Rham spectral
sequence:
$$
Q: H^{n/2}(X,\Omega^{n/2}_X)\times
H^{n/2}(X,\Omega_X^{n/2})
\xrightarrow{\cup }
H^n(X,\Omega_X^n)
\xrightarrow{\operatorname{trace}}k.
$$
\end{enumerate}
\end{theorem}

\noindent
By LPS '24, for $X = \{F=0\}
\subset \mathbb{P}^{n+1}_k$,
$Q = Q_{\operatorname{prim}} \oplus
\langle 1\rangle$ and
$$
Q_{\operatorname{prim}}:
J(X) _{(d-2)(n/2+1)}\times
J(X)_{(d-2)(n/2+1)}\xrightarrow{\cdot}
J(X)_{(d+1)(n+2)}
=\det (\text{Hess}((F)))
$$

\medskip\noindent
{\bf Non-split toric varieties.}
Recall: $X/\mathbb{C}$ is toric
if it admits an algebraic torus action
$T = (\mathbb{C}^\times)^n$, $n = \dim X$
such that there exists $x \in X$
with $T \simeq T\cdot x$ (dense open orbit).
When $X$ is normal there exists a
combinatorial object called the fan
$\Sigma = \bigcup_{k=1}^n \Sigma(k)$
consisting of cones inside $N_{\mathbb{R}}
\simeq M^4_{\mathbb{R}} \simeq \mathbb{R}^n$
where $M = \Hom(T,\mathbb{C}^\times)
\simeq \mathbb{Z}^n$
group of characters of $T$
for $\sigma \in \Sigma(k)$
if is a cone $\sigma \subseteq N_{\mathbb{R}}$
of dimension $k$, rational polyhedral
and strongly convex.

\begin{itemize}
\item for all $\sigma \in \Sigma$,
each face of $\sigma$ is in $\Sigma$,
\item for all $\sigma,\sigma' \in \Sigma$
$\sigma \cap \sigma' \in \Sigma$.
\end{itemize}

\noindent
For each $\sigma \in \Sigma$, we have
$\sigma^\vee \subseteq M_{\mathbb{R}}$
$\implies $ $\sigma^\vee \cap M \subset M$
is a semigroup, which gives
$\Spec \mathbb{C}[\sigma^\vee\cap M]
\supseteq \Spec \mathbb{C}[M]
\simeq T$.
Denote $U_\sigma = \Spec
\mathbb{C}[\sigma^\vee \cap M]$.
$\sigma \subseteq \sigma' \implies
(\sigma')^\vee \subseteq \sigma^\vee \implies 
U_\sigma \subseteq U_{\sigma'}$.

The idea is that the fan gives us the data
of how to glue the affine toric varieties
$\Spec \mathbb{C}[\sigma^\vee \cap M]$
into a toric variety.

When the rays of all cones ar $\mathbb{R}$-
linearly independent, we say
$X_{\Sigma}$ is simplicial.
\begin{itemize}
\item $X_\Sigma$ is an orbifold
(finite quotient singularity).
\item We can put homogeneous coordinates
on $X_{\Sigma}$, the {\it Cox ring}:
$$
\operatorname{Cox}(X_{\Sigma})
= \bigoplus_{\beta \in \text{Cl}(X_{\Sigma}}
H^0(\mathcal{O}_{X_{\Sigma}}(\beta)
$$
since $X_{\Sigma}\setminus T
= \bigcup_{D \in
\text{Div}_{T}(X_{\Sigma}})$.
\end{itemize}

\noindent
Each $\text{Div}_T(X_{\Sigma})
=\{D_{\rho_i}:i 1,…,r\}$,
$\Sigma(1)= \{\rho_1,…,\rho_r\}$.
$\operatorname{Cox}(X_{\Sigma})
=\mathbb{C}[X_{\rho_1},…,X_{\rho_r}]$.
This ring is naturally graded by
$\text{Cl}(X_{\Sigma})$,
$\text{deg}(X_{\rho_i})=
[D_{\rho_i}] \in \text{Cl}(X_{\Sigma})
=\bigoplus_{i=1}^r \mathbb{Z}\cdot
D_{\rho_i}/\text{div}(M)$.

A non-split torus $T$ is a $k$-linear
algebraic group such that
$T_{\overline{k}}\simeq
(\overline{k}^\times)$ is a split torus.

\begin{itemize}
\item $S^1$ is a non-split
$\mathbb{R}$-torus.
Given a $k$-torus $T$, a $T$-toric
variety is a $k$-scheme $T$ acting on $X$
such that $T_{\overline{k}}$ acts on
$X_{\overline{k}}$ becomes a usual
toric variety.
\end{itemize}

\noindent
The Galois group $\Gamma=\operatorname{Gal}
(\overline{k}/k)$ acts on $T$ and $X$.

\begin{theorem}
\label{theorem-toric-primitive-hodge}
$Y \subseteq X$, $T$-toric,
$$
H^q(Y,\Omega_Y^p)_{\operatorname{prim}}
\simeq
(Cox(X_{\overline{k}})
/\langle \frac{\partial F}{\partial x_1}
,…,\frac{\partial F}{\partial x_r}\rangle)
^\Gamma_{[\gamma]}
$$
\end{theorem}

\section{Determinantal surfaces
and Noether-Lefschetz theory}
\label{section-determinantal-surfaces-nl}

\noindent
Manuel Leal, IMPA.
GADEPS Focused Meeting June 30, 2026, IMPA.

\medskip\noindent
The talk studies determinantal surfaces
inside the parameter space of smooth
degree $d$ surfaces in $\mathbb{P}^3$.
Such a surface contains natural ACM curves,
which give extra line bundles besides
$\mathcal{O}_X(1)$.
Thus determinantal surfaces lie in the
Noether-Lefschetz locus. The main point is
that, for the degree type $(a,b)$, the
determinantal locus is itself a
Noether-Lefschetz component.

\medskip\noindent
Given a determinantal surface given
by the vanishing of the determinant
of a matrix $S$, we can define two curves:
$\mathcal{C}_+$ and $\mathcal{C}_-$
both given by the vanishing of maximal
rank minors of the matrix of $S$ after
removing everything except the highest/lowest
degree forms of the matrix.

These two curves are arithmetically
Cohen-McCaulay.

\begin{theorem}[L,LH,V, 24]
\label{theorem-llhv24}
$$
\xymatrix{
&\{(\mathcal{C}_+,X):\mathcal{C}_+ \subseteq X\}
\ar[dr]\ar[dl]\\
\mathcal{H}_{a,b}\subset
\operatorname{Hilb}_{\mathbb{P}^3}
& &
\det(a,b) \subseteq U_d
}
$$
where
$U_d$ is the open subset of
$|\mathcal{O}_{\mathbb{P}^3}(d)|$
parametrizing smooth surfaces of degree $d$.
The locus $\det(a,b)\subseteq U_d$
parametrizes surfaces cut out by the
determinant of a matrix whose entries have
degree type $(a,b)$.
The scheme $\mathcal{H}_{a,b}$ is the
Hilbert scheme component parametrizing
the corresponding arithmetically
Cohen-McCaulay curves.
\end{theorem}

\noindent
This was proved using
\begin{theorem}[Ellingsrad 75]
\label{theorem-ellin75}
$\mathcal{H}_{a,b} \subseteq
\operatorname{Hilb}_{\mathbb{P}^3}$
is an irreducible component gen. smooth
and of dimension…
\end{theorem}

\begin{proposition}
\label{proposition-}
\begin{itemize}
\item $\pi^{-1}(X)
=
\coprod_{\mathcal{C}_+ \subseteq X}
|\mathcal{O}_\ell(\mathcal{C}_+)|$.
\item $H^0(X,\mathcal{C}_+)
=
\chi(X,\mathcal{C}_+)$.
\end{itemize}
\end{proposition}

\medskip\noindent
{\bf Noether-Lefschetz theorem.} 
If $d \geq  4$ and $X \in U_d$
is very general,
$$
\text{Pic}\mathbb{P}^3
\xrightarrow{\simeq}
\text{Pic}X 
=
\langle \mathcal{O}_X(1)\rangle
\simeq \mathbb{Z}
$$
where the firs map is a resolution.

\begin{definition}
\label{definition-nl-locus}
The {\it Noether-Lefschetz locus} is
$$
NL(d)
=
\{ X \in U_d: \text{Pic}(X)\not\simeq
\mathbb{Z}\}.
$$
\end{definition}

\noindent
$$
NL(d) = \bigcup_{\operatorname{countable}}
\Sigma
$$
where $\Sigma$ are the NL-compononents.

\begin{theorem}[L,LH,V 24]
\label{theorem-llhv24ii}
$\det(a,b)$ are NL components.
\end{theorem}

\begin{theorem}[Lopez, 91]
\label{theorem-lopez-91}
A ver general degree $d$ surface
$X \supseteq \mathcal{C}_-$,
$$
\text{Pic}(X) = 
\langle \mathcal{O}_X(1),
\mathcal{O}_X(\mathcal{C}_-)\rangle.
$$
\end{theorem}

\noindent
This implies that $\det(a,b) \subseteq
NL(d)$,
and this implies that
$\det(a,b) \subseteq \Sigma$ for some
NL component.

Take $\{ X_s\} \subseteq \Sigma$
and $X_s = X \in \det(a,b)$.
There is a family of line bundles
$\{\mathcal{L}_s \in \text{Pic}(X)\setminus
\langle\mathcal{O}_{X_s}(\ell)\rangle\}$.
Since $\rho(X) = 2$ we can assume that
$\mathcal{L}_0 =
\mathcal{O}_X(\mathcal{C}_+)$.
$H^1(X,\mathcal{O}_X(\mathcal{C}_+))=0$
this implies that there exists
$\{\mathcal{C}_s \subseteq X_s\}$ such that
$\mathcal{C}_0 = \mathcal{C}_{+}$.
By the theorem of Elli, this implies that
$\mathcal{C}_s \in \mathcal{H}_{a,b}$ for
gen. s. This implies that
$X_s \supseteq \mathcal{C}_s$ is in 
$\det(a,b)$. This concludes the ideas
behind the proof.

\section{On a counterexample to a conjecture
of J. Harris for octic surfaces}
\label{section-counterexample-harris-conj}

\noindent
Hossein Movasati, IMPA.
GADEPS Focused Meeting June 30, 2026, IMPA.

\medskip\noindent
Based on works by Hossein Movasati:
2021 - Special components of
Noether-Lefschetz. 2026 - On a
counterexample…

\medskip\noindent
Let $T $ be the parameter space of smooth
degree $d$ surfaces on $\mathbb{P}^3$ over
$\mathbb{C}$.
Consider
$$
x_0^d+x_1^d+x_2^d+x_3^d
+\sum_{\text{deg}x^\alpha =d}t_\alpha
x^\alpha,
$$
where there are 165 parameters $t_\alpha$.
$d=8$.

$$
NL_\alpha
=
\{t \in T|\rho(X_t)>1\}
$$
where $\rho$ is the Picard number.
As explained in Manuel's talk,
$NL_\alpha$ is a union of enumerable
sets and
$$
5=d-3 \leq  \text{codim}(\text{components} )
\leq \binom{d-1}{3}=h^{2,0}
$$
Components which attain the upper bound
are called general components, and those
for which the right-hand inequality
is strict are called special.
General components are dense in $T$.

Harris' conjecture: the number of special
components is finite.
In the 90's, Voisin found a counter-example
for large $d$.
The conjecture is known to hold
for $d=5$. For $d=6,7$ ``the conjecture
must be true''. Thus, the next interesting
case is $d=8$.

\begin{exercise}
\label{exercise-computation-codim}
The codimension of the set of octic
surfaces containing a line is 5.
The codimension of the set of octic surfaces
containing a complete intersection
of type $(3,3)$ is 27.
Consider $x_0^8+x_1^8+x_2^8+x_3^8$
and $C_1:f_1=g_1=0$, $C_2:f_3=g_3=0$ where
the $f_i$ and $g_i$ are factors
of the Fermat equation. Then $V_{C_1}$
and $V_{C_2}$ are smooth and intersect
transversely, $\text{codim}(V_{C_1}\cap
V_{C_2})=5+27=32$.
\end{exercise}

\noindent
Let
$$
\delta_0=[C_1]+r[C_2]
\in H^2(X_0,Q) \cap H^{1,1} 
$$
where $C_1,C_2$ are as in the
exercise and $X_0$ is the Fermat,
and $r \in \mathbb{Q}$ is a parameter,
making the above a pencil of classes. 

Recall that Hodge conjecture for the case
of surfaces is called Lefschetz theorem
on $(1,1)$-classes.

Also,
$$
\delta_t \in H^2(X_t,Q),\qquad 
t \in (T,0).
$$
Local Noether-Lefschetz locus:
$$
V_{\delta_0}:=
\{t \in (T,0):
\delta_t \in H^2(X_t,\mathbb{Q})
\cap H^{1,1}(X_t)\}.
$$

In principle, this $V_{\delta_t}$
is in some component of the
Nother-Lefschetz locus.

\medskip\noindent
Conjecture: for all $r \in \mathbb{Q}$,
$r \neq  0$, except a finite number,
$V_{\delta_0}$ are distinct, smooth
varieties of codimension 31.
(Recall that $31 < 35 = h^{2,0}$.)


\section{Asymptotic behaviour of tropical
rank functions}
\label{section-asymptotic-tropical-rank}

\noindent
Eduardo Vital, Bielefeld University.
Seminar of Algebra, IMPA.
August 26, 2026, 15:30, room 228.

\medskip\noindent
{\bf Abstract.} In this talk, we study the
asymptotic behaviour of two notions of rank
for linear series on tropical curves.
Although these two notions can differ for a
fixed divisor, we show that they have the
same asymptotic growth. This leads to a
natural notion of tropical volume.

We explain the main ideas behind this
result, including the role of uniform bounds
for the Baker--Norine rank and tropical
independence. As a consequence, both notions
of rank satisfy an asymptotic
Riemann--Roch formula. We also discuss
examples where the two ranks differ at
finite level, but nevertheless have
identical asymptotic behaviour. Finally, we
relate the tropical volume to the classical
volume of divisors through tropicalization.
This is joint work with Ana Maria Botero and
Alex K\"uronya.

\medskip\noindent
{\bf Comment.}
Here I learnt that scheme theory
for tropical geometry has not been
developed. The speaker and his collaborators
are interested in developing
Cech cohomology for higher dimensions
(and in developing scheme theory
for tropical geometry).

\section{\texorpdfstring
{$C^\infty$}{C-infinity}-Superrings and
\texorpdfstring
{$C^\infty$}{C-infinity}-Superschemes}
\label{section-rizzo-smooth-superrings}

\noindent
Pedro Rizzo.
Seminar of Algebra, IMPA.
September 9, 2026, 15:30, room 228.

\medskip\noindent
{\bf Abstract.}
Superalgebraic and supergeometric structures
generalize classical nongraded structures to
graded ones. This extension introduces an
interplay between commuting and anticommuting
elements, enriching tools of commutative
algebra, differential geometry, and algebraic
geometry. Although originally motivated by
theoretical physics, supergeometry has
evolved into a robust mathematical framework
for studying objects such as supermanifolds
and superschemes.

In this talk, we introduce
$C^\infty$-superrings and their associated
geometric spaces, which we term
$C^\infty$-superschemes. We establish an
equivalence of categories between these
algebraic and geometric structures.
Furthermore, we present a generalization of
Batchelor’s theorem to $C^\infty$-superspaces
extending the classical result that every
smooth supermanifold is (noncanonically)
isomorphic to the sheaf of sections of a
vector bundle over its underlying smooth
manifold.

Conceptually, $C^\infty$-superrings
synthesize the features of $C^\infty$-rings
and superrings. Since smooth supermanifolds
form a core subclass of
$C^\infty$-superscheme, this framework
provides a natural scheme-theoretic
generalization of supergeometry in the smooth
setting. This is a joint work with Danilo
Olarte and Alexander Torres-Gómez (both at
Universidad de Antioquia) and based on the
preprints: arXiv:2508.09900v3 and
arXiv:2605.07169v2



\medskip\noindent
{\bf Introduction.}
$\mathbb{Z}_2$-structures
are inspired by commutative ``bosonic''
and anticommutative ``anti-bosonic''
elements. Its origins lie in theoretical
physics during the 70's and 80's, driven by
developments in supersymmetric field theory.

These notions are important in
supergeometry.
 
Synthetic differential geometry
also emerged in the 80's. Dubuc, Moerdijk,
Reyes and others constructed bridges between
differential and algebraic geometry.

Derived differential geometry attempts to
use techniques of algebraic geometry to
$C^\infty$ rings rather than commutative
rings.

Nishimura and Yetter explored this problem
using Weil superalgebras (which are a
particular case of $C^\infty$-superrings).
Liu and Yau introduced a notion of
$C^\infty$-superring to propose a work in
theoretical physics.

\medskip\noindent
{\bf Context and goal.}

Goal I: to establish structured mathematical
foundation of $C^\infty$ superspaces.

See results by Batchelor and Donagi-Witten.

Goal II: we propose the notion of Batchelor
space, which attempts to generalize
Batchelor's result.

\medskip\noindent
{\bf Category of $C^\infty$-rings}

\begin{definition}[C infinity ring]
\label{definition-cinf-ring}

\end{definition}

\begin{itemize}
\item A morphism s…
\item $\mathfrak{C}$ is {\it finitely
generated} if…
\item An {\it ideal} is…
\item Obvious example: the
$\mathbb{R}$-algebra of smooth functions on
a differentiable manifold.
\item Every finitely-generated
$C^\infty$-ring is isomorphic to a quotient
of $C^\infty(\mathbb{R}^n)$ by some ideal.
\item An {\it $\mathbb{R}$-point} is a
$C^\infty$-ring morphism $x:\mathfrak{C}\to
\mathbb{R}$.
\item {\it Localization} is…
\item $\mathfrak{C}$ is called {\it fair} if
it is finitely generated and
germ-determined.
\end{itemize}

\medskip\noindent
{\bf $C^\infty$-ringed space.}

\begin{definition}
\label{definition-cinf-ringed-space}

\end{definition}

\begin{itemize}
\item The stalk is localization…
\item Other expected results…
\end{itemize}

\noindent $\Spec$ and $\Gamma$ (sections)
functors are…

\begin{theorem}[Joyce]
\label{theorem-joyce-adjoint-functors}
$\Spec$ and $\Gamma$ are adjoint.
\end{theorem}

\medskip\noindent
{\bf $C^\infty$-modules and sheaves.}

\begin{itemize}
\item A {\it module over a $C^\infty$-ring} 
is…
\item 
\end{itemize}

\medskip\noindent
{\bf Superring: basic theory.}

\begin{itemize}
\item A {\it superring} is a
$\mathbb{Z}_2$-graded unitary
supercommutative ring.
\item
$\mathfrak{R}=\mathfrak{R}_{\overline{0}}
\oplus \mathfrak{R}_{\overline{1}}$
such that $\mathfrak{R}_i,
\mathfrak{R}_j \subseteq
\mathfrak{R}_{i+j}$ (mod 2).
\item A {\it morphism} is…
\item An {\it ideal} is…
\item The quotient $\mathfrak{R}/I$
is a superring…
\item The {\it canonical superideal} is…
which gives the {\it superreduced ring}…
\end{itemize}

\noindent
Batchler's result concerns precisely when
can we ``lift'' a property from the reduced
ring to all of it.

\begin{example}
\label{example-polynomial-superring}
Polynomial superring over the field
$\mathbb{K}$…
\end{example}

\begin{itemize}
\item We define $\Lambda(\mathfrak{R})$
as…
\item {\it Localization} is…
\end{itemize}

\medskip\noindent
{\bf $C^\infty$-superrings.}
Now we mix the two constructions so far.

\begin{definition}
\label{definition-cinf-superring}
A {\it $C^\infty$-superring} is…
\end{definition}



\begin{proposition}
\label{proposition-superring}
(Essentially describes the object
we shall call {\it split}.)
\end{proposition}

\noindent
[Missing content] We may thus define
super-versions of the adjoint functors
$\Spec$ and $\Gamma$…

\begin{theorem}[A]
\label{theorem-a-superring}
$s\Spec$ and $s\Gamma$ are adjoint.
\end{theorem}

\begin{theorem}[B]
\label{theorem-b-superring}
\end{theorem}

\medskip\noindent
{\bf Split spaces.}
{\it Split spaces} are… $\mathcal{X}$ is
{\it projected} if…

\begin{theorem}[C]
\label{theorem-c-superring}
Every Batchelor space is split.
\end{theorem}


\medskip\noindent
{\bf Ongoing work.}
With Bilson Castro and Alexander Torres.
Based on Bruzzo et al., understand
$C^\infty$-superstacks.

\end{document}
